\documentclass[12pt]{amsart}
\usepackage[left=20mm,right=20mm]{geometry}
\usepackage{graphicx} % Required for inserting images
\usepackage{amsmath}
\usepackage{amssymb}
\usepackage{amscd}
\usepackage{amsthm}
\usepackage{needspace}
\usepackage{enumerate}
\usepackage{xcolor}
\definecolor{linkgold}{HTML}{966e09}
\usepackage[all]{xy}
\usepackage{url}
\usepackage[colorlinks=true,allcolors=linkgold,breaklinks,pagebackref]{hyperref}
\usepackage[lite,abbrev,alphabetic]{amsrefs}
\usepackage{ragged2e}
\AddToHook{env/biblist/begin}{\RaggedRight}
\usepackage{cleveref}
\usepackage{comment}
\usepackage{todonotes}
\usepackage{tikz-cd}
\usepackage{quiver}
\usepackage{enumitem}

\makeatletter
\renewcommand{\subsection}{\@startsection{subsection}{2}%
  \z@{.7\linespacing\@plus.5\linespacing}{.35\linespacing}%
  {\normalfont\bfseries\raggedright}}
\makeatother


%theorem
\theoremstyle{definition}
\newtheorem{definition}{Definition}[section]
\newtheorem{example}[definition]{Example}

\theoremstyle{plain}
\newtheorem{theorem}[definition]{Theorem}
\newtheorem{proposition}[definition]{Proposition}
\newtheorem{lemma}[definition]{Lemma}
\newtheorem{corollary}[definition]{Corollary}
\newtheorem{conjecture}[definition]{Conjecture}
\newtheorem{hypothesis}[definition]{Hypothesis}
\newtheorem*{theorem*}{Theorem}
\newtheorem{construction}[definition]{Construction}
\newtheorem{fact}[definition]{Fact}
\newtheorem*{maintheorem}{Main theorem}


\theoremstyle{remark}
\newtheorem{remark}[definition]{Remark}
\newtheorem*{remark*}{Remark}
\newtheorem*{Ac}{Acknowledgement}
\newtheorem*{AIdisclosure}{AI disclosure}

\crefname{theorem}{Theorem}{Theorems}
\crefname{proposition}{Proposition}{Propositions}
\crefname{lemma}{Lemma}{Lemmas}
\crefname{corollary}{Corollary}{Corollaries}
\crefname{conjecture}{Conjecture}{Conjectures}
\crefname{hypothesis}{Hypothesis}{Hypotheses}
\crefname{remark}{Remark}{Remarks}
\crefname{remark*}{Remark}{Remarks}
\crefname{maintheorem}{Main theorem}{Main theorems}
\crefname{condition}{Condition}{Conditions}
\crefname{example}{Example}{Examples}

\def\N{\mathbb{N}}
\def\C{\mathbb{C}}% complex number field
\def\R{\mathbb{R}}% real number field
\def\Q{\mathbb{Q}}% rational number field
\def\Z{\mathbb{Z}}% ring of integers
\def\A{\mathbb{A}}% the ring of adeles
\def\F{\mathbb{F}}% finite field
\def\H{\mathbb{H}}% Hamilton's quaternion field




%---------------Algebraic groups------------------------
\def\GL{\mathrm{GL}}
\def\PGL{\mathrm{PGL}}
\def\SL{\mathrm{SL}}
\def\PSL{\mathrm{PSL}}
\def\O{\textnormal{O}}
\def\SO{\mathrm{SO}}
\def\GSO{\mathrm{GSO}}
\def\U{\textnormal{U}}
\def\SU{\textnormal{SU}}
\def\Sp{\mathrm{Sp}}
\def\GSp{\mathrm{GSp}}
\def\M{\mathrm{M}}% matrix space
\def\Spin{\mathrm{Spin}}
\def\PGSO{\mathrm{PGSO}}
\def\PGSp{\mathrm{PGSp}}
\def\PGO{\mathrm{PGO}}
\def\GO{\mathrm{GO}}

\def\gl{\mathfrak{gl}}
\def\sl{\mathfrak{sl}}
\def\so{\mathfrak{so}}

%--------------symbols----------------------
\def\t{\,{}^t}% transpose

\def\rd{\,\mathrm{d}}% measure

\def\idots{\reflectbox{$\ddots$}}% anti-diagonal dots

\def\Cc{C_c^{\infty}}% smooth functions

\def\ds{\displaystyle}% displaystyle

\def\bs{\backslash}% backslash
\def\rN{\mathrm{N}}

\def\tbe{\textcolor{red}{\bf To Be Edited.}}

\DeclareMathOperator{\diag}{diag} % diagonal matrix
\DeclareMathOperator{\supp}{supp}% support
\DeclareMathOperator{\ord}{ord}% order
\DeclareMathOperator{\vol}{vol}% volume
\DeclareMathOperator{\bc}{bc}
\DeclareMathOperator{\rank}{rank}% rank
\DeclareMathOperator{\Gal}{Gal}% Galois group
\DeclareMathOperator{\Res}{Res}% residue / scalar restriction
\DeclareMathOperator{\Ind}{Ind}% induction
\DeclareMathOperator{\ind}{ind}% compact induction
\DeclareMathOperator{\Hom}{Hom}% space of homomorphisms
\DeclareMathOperator{\End}{End}% space of endomorphisms
\DeclareMathOperator{\Aut}{Aut}% group of automorphisms
\DeclareMathOperator{\Lie}{Lie}% Lie algebra
\DeclareMathOperator{\Inn}{Inn}% group of inner automorphisms
\DeclareMathOperator{\rec}{rec}
\DeclareMathOperator{\Cent}{Cent}% centralizer
\DeclareMathOperator{\Irr}{Irr}% set of irreducible representations
\DeclareMathOperator{\sgn}{sgn}% the sign character
\DeclareMathOperator{\re}{Re}% real part
\DeclareMathOperator{\im}{Im}% imaginary part
\DeclareMathOperator{\Ad}{Ad}% adjoint action
\DeclareMathOperator{\ad}{ad}% adjoint action
\DeclareMathOperator{\Std}{Std}
\DeclareMathOperator{\Ker}{Ker}% kernel
\DeclareMathOperator{\Image}{Im}% image
\DeclareMathOperator{\triv}{triv}% the trivial representation
\DeclareMathOperator{\tr}{tr}% trace
\DeclareMathOperator{\id}{id}% identity operator
\DeclareMathOperator{\Sym}{Sym}% symmetrization
\DeclareMathOperator{\Int}{Int}
\DeclareMathOperator{\As}{As}
\DeclareMathOperator{\Out}{Out}
\DeclareMathOperator{\der}{der}
\DeclareMathOperator{\Ann}{Ann}
\DeclareMathOperator{\AV}{AV}


%------------fonts------------------------
%<<mathfrak>>
\newcommand{\fa}{\mathfrak{a}}
\newcommand{\fb}{\mathfrak{b}}
\newcommand{\fc}{\mathfrak{c}}
\newcommand{\fe}{\mathfrak{e}}
\newcommand{\ff}{\mathfrak{f}}
\newcommand{\fg}{\mathfrak{g}}
\newcommand{\fh}{\mathfrak{h}}
\newcommand{\fii}{\mathfrak{i}}
\newcommand{\fj}{\mathfrak{j}}
\newcommand{\fk}{\mathfrak{k}}
\newcommand{\fl}{\mathfrak{l}}
\newcommand{\fm}{\mathfrak{m}}
\newcommand{\fn}{\mathfrak{n}}
\newcommand{\fo}{\mathfrak{o}}
\newcommand{\fp}{\mathfrak{p}}
\newcommand{\fq}{\mathfrak{q}}
\newcommand{\fr}{\mathfrak{r}}
\newcommand{\fs}{\mathfrak{s}}
\newcommand{\fS}{\mathfrak{S}}
\newcommand{\ft}{\mathfrak{t}}
\newcommand{\fu}{\mathfrak{u}}
\newcommand{\fv}{\mathfrak{v}}
\newcommand{\fw}{\mathfrak{w}}
\newcommand{\fx}{\mathfrak{x}}
\newcommand{\fy}{\mathfrak{y}}
\newcommand{\fz}{\mathfrak{z}}
\newcommand{\fI}{\mathfrak{I}}
\newcommand{\fJ}{\mathfrak{J}}
\newcommand{\fK}{\mathfrak{K}}
\newcommand{\fL}{\mathfrak{L}}
\newcommand{\fP}{\mathfrak{P}}
\newcommand{\fT}{\mathfrak{T}}
\newcommand{\fW}{\mathfrak{W}}
\newcommand{\fX}{\mathfrak{X}}

%<< mathbb>>
\newcommand{\bB}{\mathbb{B}}
\newcommand{\bD}{\mathbb{D}}
\newcommand{\bE}{\mathbb{E}}
\newcommand{\bF}{\mathbb{F}}
\newcommand{\bG}{\mathbb{G}}
\newcommand{\bI}{\mathbb{I}}
\newcommand{\bJ}{\mathbb{J}}
\newcommand{\bK}{\mathbb{K}}
\newcommand{\bL}{\mathbb{L}}
\newcommand{\bM}{\mathbb{M}}
\newcommand{\bO}{\mathbb{O}}
\newcommand{\bP}{\mathbb{P}}
\newcommand{\bS}{\mathbb{S}}
\newcommand{\bT}{\mathbb{T}}
\newcommand{\bU}{\mathbb{U}}
\newcommand{\bV}{\mathbb{V}}
\newcommand{\bW}{\mathbb{W}}
\newcommand{\bX}{\mathbb{X}}
\newcommand{\bY}{\mathbb{Y}}

%<<mathcal>>
\newcommand{\cA}{\mathcal{A}}
\newcommand{\cB}{\mathcal{B}}
\newcommand{\cC}{\mathcal{C}}
\newcommand{\cD}{\mathcal{D}}
\newcommand{\cE}{\mathcal{E}}
\newcommand{\cF}{\mathcal{F}}
\newcommand{\cG}{\mathcal{G}}
\newcommand{\cH}{\mathcal{H}}
\newcommand{\cI}{\mathcal{I}}
\newcommand{\cJ}{\mathcal{J}}
\newcommand{\cK}{\mathcal{K}}
\newcommand{\cL}{\mathcal{L}}
\newcommand{\cM}{\mathcal{M}}
\newcommand{\cN}{\mathcal{N}}
\newcommand{\cO}{\mathcal{O}}
\newcommand{\cP}{\mathcal{P}}
\newcommand{\cQ}{\mathcal{Q}}
\newcommand{\cR}{\mathcal{R}}
\newcommand{\cS}{\mathcal{S}}
\newcommand{\cT}{\mathcal{T}}
\newcommand{\cU}{\mathcal{U}}
\newcommand{\cV}{\mathcal{V}}
\newcommand{\cW}{\mathcal{W}}
\newcommand{\cX}{\mathcal{X}}
\newcommand{\cY}{\mathcal{Y}}
\newcommand{\cZ}{\mathcal{Z}}

\newcommand{\bfi}{\mathbf{i}}
\newcommand{\bfj}{\mathbf{j}}
\newcommand{\bfk}{\mathbf{k}}
\newcommand{\bfI}{\mathbf{I}}
\newcommand{\bfJ}{\mathbf{J}}
\newcommand{\bfK}{\mathbf{K}}
%-------------------------------------

\numberwithin{equation}{section}
\newcommand{\disc}{\mathrm{disc}}
\newcommand{\geom}{\mathrm{geom}}
\newcommand{\spec}{\mathrm{spec}}
\newcommand{\elliptic}{\mathrm{ell}}
\newcommand{\scusp}{\mathrm{scusp}}
\newcommand{\cusp}{\mathrm{cusp}}
\newcommand{\semi}{\mathrm{ss}}
\newcommand{\reg}{\mathrm{reg}}
\newcommand{\temp}{\mathrm{temp}}
\newcommand{\gen}{\mathrm{gen}}
\newcommand{\bdd}{\mathrm{bdd}}
\newcommand{\unit}{\mathrm{unit}}
\newcommand{\spl}{\mathrm{spl}}
\newcommand{\red}{\mathrm{red}}
\newcommand{\pl}{\mathrm{pl}}
\newcommand{\rG}{\mathrm{G}}
\newcommand{\Art}{\mathrm{Art}}
\newcommand{\ab}{\mathrm{ab}}



\newcommand{\Jor}[6]
{
    \begin{pmatrix}
    #1               &        #4            &   #5^* \\
    #4^*             &        #2            &   #6 \\
    #5               &        #6^*          &   #3
    \end{pmatrix}
}

\newcommand{\unimat}[3]
{
    \begin{pmatrix}
    1              &        #1           &   #3 \\
    0              &        1            &   #2 \\
    0              &        0            &   1
    \end{pmatrix}
}

\newcommand{\opunimat}[3]
{
    \begin{pmatrix}
    1              &        0            &   0 \\
    #1             &        1            &   0 \\
    #3             &        #2           &   1
    \end{pmatrix}
}


\newcommand{\esixroot}[6]{
  \begin{array}{ccccc}
     &   &  #2 &   &   \\
  #1 & #3 & #4 & #5 & #6 \\
  \end{array}
}

\newcommand{\Temp}{\mathrm{Temp}}
\newcommand{\abs}[1]{\lvert #1 \rvert}
\newcommand{\heart}{\ensuremath\heartsuit}
\newcommand{\spade}{\ensuremath\spadesuit}
\newcommand{\st}{\, \vert \,}
\newcommand{\pr}{\mathrm{pr}}
\newcommand{\GSpin}{\mathrm{GSpin}}
\newcommand{\Pin}{\mathrm{Pin}}
\newcommand{\Isom}{\mathrm{Isom}}
\newcommand{\Norm}{\mathrm{Norm}}
\newcommand{\dist}{\mathrm{dist}}
\newcommand{\Trans}{\mathrm{Trans}}
\newcommand{\St}{\mathrm{St}}
\newcommand{\sn}{\mathrm{sn}}



\DefineSimpleKey{bib}{arxiv}
\newcommand\myurl[1]{\url{#1}}
\newcommand\arxiv[1]{available at \href{https://arxiv.org/abs/#1}{arXiv:#1}}
\def\MR#1{\quad \href{http://www.ams.org/mathscinet-getitem?mr=#1}{MR#1}}
\BibSpec{article}{%
  +{}  {\PrintAuthors}                {author}
  +{,} { \textit}                     {title}
  +{.} { }                            {part}
  +{:} {\textit}                     {subtitle}
  +{,} {\PrintContributions}         {contribution}
  +{,} {\PrintConference}            {conference}
  +{}  {\PrintBook}                   {book}
  +{,} { }                            {booktitle}
  +{,} { }                            {journal}
  +{,} { \textbf}                     {volume}
  +{,} { \issuetext}                  {number}
  +{,} { \PrintDateB}                 {date}
  +{,} { pp.~}                       {pages}
  +{,} { available at \eprint}        {eprint}
  +{;} { \arxiv}               {arxiv}
  +{.\!\!}  {\SentenceSpace \PrintReviews} {review}
  +{,} { \PrintDOI}                   {doi}
}




\allowdisplaybreaks[1]
\emergencystretch=2em
\hypersetup{
  bookmarksdepth=2,bookmarksnumbered=true,bookmarksopen=false,
  pdftitle={Arthur packets for G2 over p-adic fields},
  pdfauthor={Yugo Takanashi}
}

\title[Arthur packets for $\rG_2$ over $p$-adic fields]{Arthur packets for $\rG_2$ over $p$-adic fields}
\author{Yugo Takanashi}
\date{}

\begin{document}
\begin{abstract}
We describe Arthur packets for the split reductive group $\rG_2$ over a $p$-adic field and compute their stable and endoscopic character distributions.
The construction uses the twisted Aubert duality recently studied by Cheng--Kaletha in \cite{ChengKaletha2026-AubertBernsteinInvolutionsDisconnectedGroups}.
We identify the corresponding Xu packets on $\PGSO_8$ and their triality-twisted characters. 
We also extend the result to long-root and short-root Arthur packets corresponding to unitary generic representations of $\PGL_2$.
\end{abstract}

\maketitle
\setcounter{tocdepth}{1}
\tableofcontents

\section*{Introduction}


Let $F$ be a finite extension of $\Q_p$. 
Gan and Savin constructed the local Langlands correspondence for the split group $\rG_2$ over $F$ \cite{GanSavin2023-LocalLanglandsConjectureG2}. 
The author proved the ordinary endoscopic character relations for its tempered packets and the triality-twisted character relations for their lifts to $\PGSO_8$ \cite{TakanashiEndoscopyLocal}. 
Several families of local non-tempered Arthur packets were previously constructed using theta correspondences:
the cubic unipotent packets arose from the work of Huang--Magaard--Savin \cite{HuangMagaardSavin1998-UnipotentRepresentationsG2} and
Gan--Gurevich--Jiang \cite{GanGurevichJiang2002-CubicUnipotentArthurParameters}, and the short-root packets were constructed through the Rallis--Schiffmann correspondence by Gan--Gurevich \cite{GanGurevich2006-NontemperedAPacketsG2}.

The goal of this article is to construct the non-tempered Arthur
packets of $\rG_2$ from its tempered packets and to prove both ordinary endoscopic character identities for $\rG_2(F)$ and triality-twisted endoscopic character relations for $\PGSO_8(F)$.

\begin{maintheorem}[Informal; see Theorem~\ref{thm:main-ecr}]\label{thm:main-intro}
For every non-tempered Arthur parameter $\psi$ of $G_2(F)$, we construct a finite
packet $\Pi_\psi$ consisting of irreducible representations of $G_2(F)$ with a surjective map
\begin{align*}
      \Pi_{\psi} \twoheadrightarrow \Irr_{\C}(\cS_{\psi}),
\end{align*}
and a stable virtual character $S\Theta_\psi$. These packets have the following properties:
\begin{itemize}
      
\item They satisfy the endoscopic character identities for the endoscopic groups of $G_2(F)$ and the triality-twisted endoscopic character identities for $\PGSO_8(F)$.

\item The triality-twisted transfer of $S\Theta_\psi$ to $\PGSO_8$
is an explicit signed sum of canonical twisted characters in the
corresponding distinguished Xu packet constructed by Xu in \cite{Xu2024-ArthurPacketsQuasisplitGSp2nGO2n}.
This property characterizes the local Arthur packets of $G_2(F)$ uniquely.

\item For non-tempered co-tempered parameters, an explicit sign
formula relates their stable characters to the Aubert duals of the corresponding tempered stable characters.
For co-tempered parameters, the map above is also bijective, and the compatibility with Aubert duality also gives a characterization of the packets.

\end{itemize}
\end{maintheorem}



For co-tempered parameters, we apply Aubert duality to the tempered packet obtained by interchanging the Deligne and Arthur $\SL_2$ factors.
The mixed cases also use Aubert duality to reduce the problem to long-root parameters and compute the packet explicitly by using known computations on composition series. 
For $\rG_2$, these calculations are facilitated by Qin's computations of Aubert duality
\cite{Qin2026-AubertZelevinskyInvolutionG2HeckeAlgebras}, together with the composition series of Mui\'c \cite{Muic1997-UnitaryDualG2}.
Our use of Aubert duality is inspired by Arthur's treatment
\cite{Arthur2013-EndoscopicClassificationRepresentations}*{Section 7.1}, \cite{Hiraga2004-FunctorialityZelevinskiInvolutions}
and the work of Atobe--Gan--Ichino--Kaletha--M\'inguez--Shin
\cite{AtobeGanEtAl2025-LocalIntertwiningRelationsCoTemperedPackets}.

For archimedean fields, the short root Arthur packets are constructed by Gan--Gurevich \cite{GanGurevich2006-NontemperedAPacketsG2}*{Sections 10 and 11} and the subregular packets are described in  Gan--Gurevich--Jiang \cite{GanGurevichJiang2002-CubicUnipotentArthurParameters}*{Section 6} via theta correspondences.
We will treat the archimedean case in a forthcoming paper, using an approach similar to that of this article.

The compatibility of our construction of Arthur packets with the Huang--Magaard--Savin, Gan--Gurevich--Jiang, and Gan--Gurevich constructions will be established in a forthcoming paper, as a step towards the endoscopic classification of discrete automorphic representations of $\rG_2$.
In recent work, Cunningham--Hazeltine--Liu--Lo--Ray--Xu--Sch\"urmann develop a framework to study the Kottwitz--Shelstad transfer for ABV-packets. 
They also suggest applying their framework to the triality-twisted endoscopy relating $\rG_2$ and $\PGSO_8$; see \cite{CunninghamEtAl2026-GeometrizationKottwitzShelstadTransfer}*{Section 1.5}.
It may be interesting to compare our packets with the corresponding ABV-packets.


\subsection*{Outline of the paper}
\begin{itemize}
      \item Section~\ref{sec:generalities} collects the general representation-theoretic notation, parameter conventions, transfer normalizations, and general results on Aubert duality. 
      \item Sections~\ref{sec:preliminaries-G2-PGSO8}--\ref{sec:twisted-endoscopy} specialize these inputs to
      $\rG_2$ and the orthogonal groups, including the comparison with Qin's notation, the tempered packet identities, and the signs for triality transfer. 
      \item Sections~\ref{sec:principal-parameter}--\ref{sec:mixed-short} compute the elements of A-packets and the character identities for each family of parameters.
      \item Section~\ref{sec:summary} states and proves the main theorem from these calculations. 
      \item  Section~\ref{sec:A-parameter-unbounded-ecr} also treats root parameters arising from unitary generic representations of $\PGL_2(F)$, including complementary series.
\end{itemize}

  
\begin{Ac}
      The author thanks Sug Woo Shin for his suggestion to follow the approach of Atobe--Gan--Ichino--Kaletha--M\'inguez--Shin \cite{AtobeGanEtAl2025-LocalIntertwiningRelationsCoTemperedPackets} and for his careful comments on the draft.
      This work was supported by JSPS KAKENHI Grant Number JP26KJ0167. 
\end{Ac}

\smallskip
\noindent\textbf{AI disclosure.}
The author first wrote the initial draft without using AI and subsequently used AI to polish the manuscript. The author takes full responsibility for the content of this paper.





\section{Generalities}\label{sec:generalities}

\subsection{Basic conventions}\label{subsec:basic-conventions}

Throughout, $F$ is a finite extension of $\Q_p$. 
We fix a nontrivial additive character $\psi_F \colon F \to \C^\times$.
In this section, let $G$ denote a connected reductive group over $F$ and $P = MN$ a parabolic subgroup of $G$ with Levi factor $M$ and unipotent radical $N$.
This section establishes general conventions and tools used below.

We use Artin reciprocity $\Art_F \colon F^{\times} \to W_F^{\ab}$, which sends a uniformizer to an arithmetic Frobenius element.
A character $\chi$ of $F^{\times}$ is viewed as a character of $W_F$ via $\chi \circ \Art_F^{-1}$.
In particular, $\nu(w) = \abs{w}_F$.
We set $\nu_n=|\det|_F$ on $\GL_n(F)$.
For the group $\GSO_{n}$, let $\mathrm{sim} \colon \GSO_n(F) \to F^{\times}$ denote the similitude character.

\subsection{Representations and normalized inductions}\label{subsec:general-representations}

For a connected reductive $F$-group $G$, write $\Irr(G(F))$ for
the isomorphism classes of irreducible admissible complex
representations, and $\operatorname{JH}$ for irreducible constituents
of a finite-length representation. Let $R(G)$ be the
Grothendieck group of finite-length smooth complex representations
of $G(F)$, and write
$\Theta_\pi(f)=\tr\pi(f)$ for $f\in C_c^\infty(G(F))$.
If a representation $\pi$ has a central character, let $\omega_\pi$ denote it.

Fix a minimal $F$-parabolic $P_{0, G}=M_{0, G}N_{0, G}$.
Standard parabolics contain $P_{0, G}$, and their standard
Levi factors contain $M_{0, G}$.
We sometimes omit the subscript $G$ when the group is clear from context.
For a parabolic $P=MN$ of $G$, let $\delta_P$ be its modulus
character and put
\[
 I_P^G(\sigma)=
 \operatorname{Ind}_{P(F)}^{G(F)}(\delta_P^{1/2}\sigma),
 \qquad
 r_P^G(\pi)=\delta_P^{-1/2}\otimes\pi_{N(F)}.
\]
Write $\overline P=M\overline N$ for the parabolic opposite to $P$.
Here $\sigma$ is extended trivially across $N(F)$, and
$\pi_{N(F)}$ denotes the unnormalized Jacquet module.
When a standard Levi is displayed in an induction or quotient
subscript, it denotes the standard parabolic with that Levi.
For orthogonal groups, $\tau\rtimes\rho$ denotes normalized
induction from the displayed Levi $\GL_n\times\SO_m$, with
$\tau$ on $\GL_n(F)$ and $\rho$ on $\SO_m(F)$.


Let $X_F^*(M)$ be the group of $F$-rational algebraic characters
of $M$, and put
\[
 \mathfrak a_M=\Hom_{\mathbf Z}(X_F^*(M),\mathbf R).
\]
Write $\mathfrak a_M^*$ for its real dual and
$\mathfrak a_{M,\C}^*=\mathfrak a_M^*\otimes_{\mathbf R}\C$.
Define $H_M:M(F)\to\mathfrak a_M$ by
$\langle\xi,H_M(m)\rangle=\log|\xi(m)|_F$.
For $\lambda\in\mathfrak a_{M,\C}^*$, set
\[
 \nu_M^\lambda(m)=\exp\langle\lambda,H_M(m)\rangle.
\]
We also write $\nu^\lambda$ when $M$ is understood.
Thus $\nu_M^{s\xi}=|\xi|_F^s$ for $\xi\in X_F^*(M)$.
If $\sigma$ is irreducible tempered and
$\lambda\in\mathfrak a_M^*$ is strictly positive relative to a
standard $P$, then $I_P^G(\sigma\otimes\nu_M^\lambda)$ is a
Langlands standard module. Its unique irreducible quotient is
denoted by $J_P^G(\sigma\otimes\nu_M^\lambda)$.
The notation $J_P^G(\eta)$ is used only when the displayed
induction is such a standard module. We also allow $P=G$ and
$\lambda=0$, giving an irreducible tempered standard module.

We now specialize to $G = \GL_2$.
Let $B_2$ be the upper triangular Borel of $\GL_2$. For arbitrary
smooth characters $a,b$ of $F^\times$, put
\[
 \pi(a,b)=I_{B_2}^{\GL_2}(a\otimes b).
\]
Here $(a\otimes b)(\diag(x,y))=a(x)b(y)$.
This is the entire normalized principal series, which may be
reducible; if $a,b$ are unitary, it is irreducible tempered.
Write $\St_G$ for the Steinberg representation of $G(F)$,
$\St_n=\St_{\GL_n}$, and, for a smooth character $\chi$ of $F^\times$,
$\St_2(\chi)=\St_2\otimes(\chi\circ\det)$, so
$\St_2=\St_2(1)$. A twist $\tau\nu_2^s$ means
$\tau\otimes|\det|_F^s$.

We shall also use the standard $\GL_2$ exact sequence
\begin{equation}\label{eq:GL2-length-two}
 0\longrightarrow\St_2(\chi)\nu_2^{1/2}
 \longrightarrow\pi(\chi\nu,\chi)
 \longrightarrow(\chi\circ\det)\nu_2^{1/2}\longrightarrow0,
\end{equation}
valid for every smooth character $\chi$ of $F^\times$; see
\cite{GanSavin2023-HoweDualityDichotomyExceptionalThetaCorrespondences}*{Section 3.1}.

\subsection{Parameters and infinitesimal characters}\label{subsec:parameters-and-infinitesimal-characters}

The parameter conventions needed here concern split connected reductive
groups $G$. We fix a pinned dual group $\widehat G$, write
${}^LG=\widehat G\times W_F$, and display only the $\widehat G$-component
of each parameter; its $W_F$-component is the projection onto $W_F$.
Write $L_F=W_F\times\SL_2^D(\C)$, where $D$ denotes the Deligne factor.
A Langlands parameter is an admissible homomorphism
$\varphi:L_F\to\widehat G$, algebraic on $\SL_2^D$; it is tempered
when its restriction to $W_F$ is bounded. Parameters are considered
up to $\widehat G$-conjugacy. Put
$d_w=\diag(|w|_F^{1/2},|w|_F^{-1/2})$.

\begin{definition}\label{def:parameters}
An Arthur parameter for $G$ is an admissible homomorphism
\[
 \psi:W_F\times\SL_2^D(\C)\times\SL_2^A(\C)\longrightarrow\widehat G
\]
that is bounded on $W_F$ and algebraic on the two $\SL_2$ factors;
$A$ denotes the Arthur factor. Its associated Langlands parameter
and its swapped parameter are
\begin{align}
 \phi_\psi(w,x)&=\psi(w,x,d_w),\label{eq:associated}\\
 \widehat\psi(w,x,y)&=\psi(w,y,x).\notag
\end{align}
We call $\psi$ tempered if it is trivial on $\SL_2^A$, and
co-tempered if it is trivial on $\SL_2^D$; equivalently,
$\widehat\psi$ is tempered.

We also set 
\begin{align*}
      s_{\psi}=\psi(1, I_2, -I_2).
\end{align*}
\end{definition}

Let $S_a=\Sym^{a-1}(\C^2)$. In any fixed linear representation
of $\widehat G$, a summand $\rho\boxtimes S_a^D\boxtimes S_b^A$
of $\psi$ contributes
\begin{equation}\label{eq:parameter-block}
 \bigoplus_{j=0}^{b-1}\rho\nu^{(b-1)/2-j}\boxtimes S_a^D
\end{equation}
to $\phi_\psi$. Indeed, the diagonal torus acts on $S_b$ with weights
$b-1,b-3,\ldots,1-b$.

The infinitesimal character of a Langlands parameter $\varphi$ is the
$\widehat G$-conjugacy class of the homomorphism
\[
 \lambda_\varphi(w)=\varphi(w,d_w).
\]
For an Arthur parameter this gives
\begin{equation}\label{eq:infinitesimal-Arthur}
 \lambda_\psi(w)=\lambda_{\phi_\psi}(w)=\psi(w,d_w,d_w).
\end{equation}
In particular, swapping the two $\SL_2$ factors preserves the
infinitesimal character. 

The notion of infinitesimal character is compatible with parabolic induction and Jacquet module functors for quasi-split classical groups or similitude groups, see \cite{Xu2024-ArthurPacketsQuasisplitGSp2nGO2n}*{Corollary 5.5}.


\subsection{Endoscopic transfer, induction, and stability}\label{subsec:transfer-normalization}

Let $G$ be quasi-split, fix a pinning and its Whittaker datum
using $\psi_F$, and let $a$ be a finite-order $F$-automorphism
preserving the pinning. 
Let $(H, s, \cH, \xi)$ be an endoscopic group of $(G, a)$.
In this paper, we only consider the case $$\cH=\widehat H\rtimes W_F$$ and $\xi$ is the natural embedding $\cH\hookrightarrow{}^LG$.


We distinguish the coset $G(F) \rtimes a$
from the disconnected group $G(F)\rtimes\langle a \rangle$.
For transfer on a fixed coset $G(F) \rtimes a$, we use the corrected
Whittaker-normalized Artin factor $\Delta_A$ of
\cite{AtobeGanEtAl2025-LocalIntertwiningRelationsCoTemperedPackets}*{Appendix F.7},
denoted $\Delta'_\lambda$ in
\cite{KottwitzShelstad2012-SplittingInvariantsSignConventions}*{equation (5.5.2)}.
We use the discriminant-normalized orbital integrals and compatible measures of
\cite{MoeglinWaldspurger2016-StabilisationFormuleTracesTordue1}*{I.2.4}.

Here $S^H(\delta,f^H)$ denotes a stable orbital integral and
$I^{G \rtimes a}(\gamma,f)$ a twisted orbital integral.
Conjugacy on $G(F) \rtimes a$ is under $G(F)$, as in
\cite{Kaletha2022-LocalLanglandsConjecturesDisconnectedGroups}*{Lemma 4.6}.
Put
\[
 \Delta(\delta,\gamma)
       =\Delta_A(\delta,\gamma)/\Delta_{\mathrm{IV}}(\delta,\gamma).
\]
Following \cite{TakanashiEndoscopyLocal}*{Section 1.3}, matching functions
satisfy, for strongly $G$-regular $\delta$,
\[
 S^H(\delta,f^H)=d(a)^{1/2}
 \sum_{\gamma\leftrightarrow\delta}
 \frac{\Delta(\delta,\gamma)}{[Z_G(\gamma)(F):G_\gamma(F)]}
 I^{G \rtimes a}(\gamma,f).
\]
The sum is over rational conjugacy classes matching $\delta$, and
$G_\gamma=Z_G(\gamma)^\circ$ is the connected centralizer used in
$I^{G \rtimes a}$. For the Lie algebra $\mathfrak t_G$ of the pinned
maximal torus $T_G$, put
\[
 d(a)=
 \left|\det(1-a\mid\mathfrak t_G/\mathfrak t_G^a)\right|_F.
\]
Thus $d(1)=1$; for pinned triality of type $D_4$, $d(a)=|3|_F$.

Write $\operatorname{Tran}_H^{G \rtimes a}$ for spectral transfer from
stable distributions on an endoscopic group $H$ to invariant distributions
on $G(F) \rtimes a$. Thus
\[
 \operatorname{Tran}_H^{G \rtimes a}(S)=T
 \quad\text{means}\quad T(f)=S(f^H).
\]
For ordinary endoscopy take $a=1$ and omit it from the
notation. A stable distribution is one that factors through stable
orbital integrals. In particular, stability of a character is tested
on stable conjugacy classes of regular semisimple elements.



We specify the Levi data used for induction. Let
$\eta:{}^LH\hookrightarrow{}^LG$ be the fixed endoscopic embedding,
and write $\dot s\in\widehat G\rtimes\widehat a$ for its
endoscopic element. Assume the ambient datum is elliptic and use
embedded data with
$\eta(\widehat H)=\Cent(\dot s,\widehat G)^\circ$ and
$\eta({}^LH)\subset\Cent(\dot s,{}^LG)$, as in the applications.
For an $F$-parabolic $P_H=M_HN_H$, put
\[
 \widehat A_{M_H}=(Z(\widehat M_H)^{W_F})^\circ,
 \qquad \mathcal M=\Cent(\eta(\widehat A_{M_H}),{}^LG).
\]
The Weil action here is supplied by ${}^LM_H$.
By \cite{MoeglinWaldspurger2016-StabilisationFormuleTracesTordue1}*{I.3.4},
after conjugation $\mathcal M={}^LM$ for an $a$-stable Levi
$M$ of $G$. Choose $\lambda\in X_*(\widehat A_{M_H})$ in the open
cone defining $\widehat P_H$, avoiding the kernels of the nonzero
restrictions of $\widehat G$-roots to $\eta(\widehat A_{M_H})$.
Define $\widehat P$ by the nonnegative weights of $\eta\circ\lambda$,
and let $P=MN$ be the corresponding parabolic. Restriction gives
\[
 \eta_M:{}^LM_H\hookrightarrow{}^LM,
 \qquad
 \eta_M(\widehat M_H)=\Cent(\dot s_M,\widehat M)^\circ,
 \qquad \dot s_M\in\widehat M\rtimes\widehat a.
\]
Here $\dot s_M$ is the conjugate of $\dot s$ in the chosen Levi.
We assume this inherited datum is relevant, so matching strongly
regular $F$-points exist. Its auxiliary data, transfer factor, and
measures are inherited from the ambient datum as in loc.~cit.
We call this the \emph{inherited Levi datum} and denote its spectral
transfer by $\operatorname{Tran}_{M_H,\eta_M}^{M \rtimes a}$.




\begin{lemma}\label{lem:induction-transfer}
For the inherited Levi datum above and a stable virtual character
$S$ of $M_H(F)$, one has
\[
 \operatorname{Tran}_{H,\eta}^{G \rtimes a}(I_{P_H}^H S)
   = I_P^G\bigl(\operatorname{Tran}_{M_H,\eta_M}^{M \rtimes a}(S)\bigr).
\]
Induction of distributions is defined by normalized constant terms;
on the right it is evaluated on $G(F)\rtimes a$.
For ordinary induction, $I_P^G S_M$ is stable whenever $S_M$ is a
stable virtual character of $M(F)$.
\end{lemma}
\begin{proof}
If $f^H$ is a transfer of $f$, write $f_P$ and $(f^H)_{P_H}$ for
normalized constant terms. Descent with the inherited factors and
measures gives
\[
 [(f^H)_{P_H}]=[(f_P)^{M_H}]
\]
in the quotient of test functions on $M_H(F)$ by the common kernel
of stable orbital integrals. This is
\cite{MoeglinWaldspurger2016-StabilisationFormuleTracesTordue1}*{
I.4.11, condition (2)}, with the data of I.3.4.
Evaluation on $S$ and the constant-term character formula prove the
identity. For the last assertion, normalized constant terms induce a
map from the stable quotient for $G$ to that for $M$; see loc.~cit.,
I.3.1. Its dual sends $S_M$ to the stable character $I_P^G S_M$.
\end{proof}

For the standard nonelliptic Levi datum we also use the following
ordinary transfer identity. Let $G$ be split, let $P=MN$ be standard,
and take $H=M$ with the standard inclusion ${}^LM\hookrightarrow{}^LG$.
Choose $\dot s\in Z(\widehat M)^\circ$ with
$\Cent(\dot s,\widehat G)^\circ=\widehat M$, and inherit the
Whittaker normalization, auxiliary data, and measures from $G$.
Then, for every stable virtual character $S$ of $M(F)$,
\[
 \operatorname{Tran}_{M}^{G}(S)=I_P^G(S).
\]
This is the identity in
\cite{Hiraga2004-FunctorialityZelevinskiInvolutions}*{Section 6,
Lemma 6.1, p.~1642}, with the datum on $M$ equal to the identity
datum. Equivalently, normalized constant-term descent followed by
evaluation on $S$ gives normalized induction. This separate statement
applies to the nonelliptic Levi datum; the ambient ellipticity
assumption preceding Lemma~\ref{lem:induction-transfer} is not used.

For later use, we record a general criterion for automatic stability
of Levi characters.

\begin{lemma}\label{lem:central-quotient-stability}
Let $L$ be a product of general linear groups over $F$, and let
$Z\subset L$ be a central split torus. Every finite-length character
of $(L/Z)(F)$ is stable.
\end{lemma}
\begin{proof}
If $T\subset L/Z$ is a maximal torus, its inverse image $T_L$ in $L$
is a product of multiplicative groups of finite etale $F$-algebras.
Hence $H^1(F,T_L)=H^1(F,L)=1$. 
The central-extension cohomology
sequences show that
\[
 \ker\bigl(H^1(F,T)\longrightarrow H^1(F,L/Z)\bigr)=1,
\]
and the connecting map from $H^1(F,T)$ to $H^2(F,Z)$ is injective,
and is compatible with the connecting map for $L/Z$.
Thus a strongly regular semisimple stable conjugacy class in $L/Z$ contains
only one rational conjugacy class. Every invariant character is
therefore stable.
\end{proof}

\subsection{Canonical extensions for an $S_3$-action}\label{subsec:canonical-extensions-s3-action}

Let $G$ be a connected reductive group over $F$ with an action of $S_3$ by $F$-automorphisms.
Write $A_3 \subset S_3$ for the alternating subgroup and choose a generator $a$ of $A_3$.


\begin{proposition}\label{prop:canonical}
      If $\pi$ is an admissible representation of $G(F)$ such that $\End_{G}(\pi)=\mathbb{C}$ and $\pi \cong \pi^g$ for $g \in S_3$, then there exists an extension $\pi^{+}$ of $\pi$ to $G(F) \rtimes S_3$.
      The restriction of $\pi^{+}$ to $G(F) \rtimes A_3$ is uniquely determined by $\pi$ up to isomorphism.
\end{proposition}

\begin{proof}
      Choose operators $U, V$ implementing generators of orders $3$ and $2$ of $S_3$. 
      We may assume that $U^3 = 1$ and $V^2 = 1$.
      Schur's lemma implies that 
      \begin{align*}
            VUV^{-1} = c U^{-1}
      \end{align*}
      for some $c \in \mathbb{C}^{\times}$.
      Then, we have $c^3 = 1$.
      By replacing $U$ by $cU$, we have $VUV^{-1} = U^{-1}$.
      Thus, we obtain a representation of $G(F) \rtimes S_3$ on the same space as $\pi$.
      Both the trivial and sign characters of $S_3$ are trivial on $A_3$, so the restriction of $\pi^{+}$ to $G(F) \rtimes A_3$ is uniquely determined by $\pi$ up to isomorphism.
\end{proof}

For a finite-length representation $V$ of $G(F) \rtimes S_3$, we use the twisted character on $G(F) \times a$ defined by
\begin{align*}
      \Theta_V(f) = \tr \left( \int_{G(F)} V(g \rtimes a) f(g) dg \right)
\end{align*}
for $f \in C_c^{\infty}(G(F))$.

\begin{corollary}
      If $\pi$ is an admissible representation of $G(F) \rtimes S_3$ of length $2$ as a $G(F)$ representation with an irreducible socle $S$ and cosocle $C$ as $G(F)$ representations, then we have  
      \begin{align*}
            \Theta_{\pi}(f)|_{G(F) \times a} = \Theta_{\widetilde{S}}(f) + \Theta_{\widetilde{C}}(f)
      \end{align*}
\end{corollary}

\subsection{Aubert involution and endoscopic transfer}\label{subsec:aubert-involution-and-endoscopic-transfer}

For a connected reductive $F$-group $G$ put
\begin{equation}\label{eq:Aubert-operator}
 \mathcal D_G=\sum_{P=MN\supset P_{0,G}}
    (-1)^{\dim(A_M/A_G)}I_P^G r_P^G.
\end{equation}
Here $A_M$ is the maximal split torus in the center of $M$.
For irreducible $\pi$ with supercuspidal support $(L,\sigma)$, put
$\epsilon_G(\pi)=(-1)^{\dim(A_L/A_G)}$.

\begin{proposition}\label{prop:Aubert}
There is a positive irreducible Aubert dual $\widehat\pi$ such that
\[
 \mathcal D_G[\pi]=\epsilon_G(\pi)[\widehat\pi].
\]
On the subcategory whose irreducible subquotients have supercuspidal
support on Levis of a fixed relative semisimple rank, the positive dual extends
to an exact covariant functor. It is involutive on isomorphism classes
and satisfies
\[
 \widehat{I_P^G(\sigma)}\simeq I_{\overline P}^G(\widehat\sigma).
\]
For supercuspidal $\pi$, one has $\widehat\pi=\pi$.
\end{proposition}
\begin{proof}
Let $r$ be the relative semisimple rank of $G$ and let $t$ be that
of a Levi in the supercuspidal support. Give the term indexed by a subset $J$ of the relative simple roots
degree $|J|$ in the Aubert complex $C_\pi$. Since the exponent in
\eqref{eq:Aubert-operator} is $r-|J|$, one has
$\mathcal D_G[\pi]=(-1)^r\chi(C_\pi)$.
The cohomology is concentrated in the support degree $t$ and equals
the positive dual. Hence
\[
 \mathcal D_G[\pi]=(-1)^{r+t}[\widehat\pi]
      =(-1)^{r-t}[\widehat\pi]
      =\epsilon_G(\pi)[\widehat\pi].
\]
We use Aubert's irreducibility and involutivity theorems and the
compatibility with induction in the formulation recalled in
\cite{AtobeGanEtAl2025-LocalIntertwiningRelationsCoTemperedPackets}*{
Theorems B.2.1 and B.2.3}. For exactness, see also the homological
duality of Schneider--Stuhler
\cite{SchneiderStuhler1997-RepresentationTheorySheaves}*{Section III.3}.
For a supercuspidal representation every proper
Jacquet module is zero, so only $P=G$ remains in
\eqref{eq:Aubert-operator} and the sign is $+1$.
\end{proof}

We write $D_A$ for this positive Aubert functor, on the subcategories of fixed relative support-Levi rank in Proposition~\ref{prop:Aubert}.

For a split semisimple group $G$ and a finite-order
pinning-preserving automorphism $a$, we also use the twisted
operator
\begin{equation}\label{eq:twisted-Aubert-operator}
 \mathcal D_G^a=
  \sum_{\substack{P=MN\supset P_{0,G}\\a(P)=P}}
    (-1)^{\dim A_M^a}I_P^G r_P^G.
\end{equation}
Here $P_{0,G}$ is the pinned Borel. The exponent is the dimension
of the identity component of the $a$-fixed subgroup of $A_M$,
and the induced operators are evaluated on $G(F)\rtimes a$.
Since $G$ is semisimple,
$\mathcal D_G^1=\mathcal D_G$.

Compatibility with ordinary endoscopic transfer is due to Hiraga
\cite{Hiraga2004-FunctorialityZelevinskiInvolutions}*{Theorems 6.6--6.7}.
In our convention, for a split endoscopic group $H$
of a split $G$ it reads
\[
 \mathcal D_G\operatorname{Tran}_H^G
   =(-1)^{\dim A_H-\dim A_G}
       \operatorname{Tran}_H^G\mathcal D_H.
\]
Indeed, Hiraga's operator is $(-1)^{r(G)}\mathcal D_G$, where
$r(G)$ is the semisimple split rank, and his transfer sign is
$(-1)^{a(G)-a(H)}$, where $a(G)=r(G)+\dim A_G$.
In particular, the displayed sign is $+1$ for split semisimple
groups. 
We use Xu's extension to twisted endoscopy
\cite{Xu2017-MoeglinParametrizationArthurPackets}*{Appendix A}.

\begin{proposition}\label{prop:Xu-transfer}
Let $H$ be an endoscopic group of $(G, a)$, where the groups
are split semisimple and $a$ preserves a pinning. Assume the
datum is realized by an $L$-embedding as in
\cite{Xu2017-MoeglinParametrizationArthurPackets}*{Appendix A}, with
compatible transfer normalizations. Then, on stable virtual characters,
\begin{equation}\label{eq:Xu-transfer}
 \mathcal D_G^a\operatorname{Tran}_H^{G \rtimes a}
   =\operatorname{Tran}_H^{G \rtimes a}\mathcal D_H.
\end{equation}
The Aubert operator on the right preserves stable virtual characters.
\end{proposition}

\begin{proof}
We give the indexing in Xu's Jacquet-transfer identity.
Conjugate the embedding $\eta$ so that
$\eta(\widehat T_H)=(\widehat T_G^{\widehat a})^\circ$
and $\eta(\widehat B_H)\subset\widehat B_G$, and identify ${}^LH$
with its image. For an $a$-stable standard $P=MN$, put
\[
 W_{G, a}=W(\widehat G,\widehat T_G)^{\widehat a},
 \quad W_{M,a}=W(\widehat M,\widehat T_G)^{\widehat a},
 \quad W_H=W(\widehat H,\widehat T_H)\subset W_{G, a}.
\]
Choose a torus $A\subset(\widehat T_G^{\widehat a})^\circ$
with ${}^LM=\Cent(A,{}^LG)$, and define
\[
 \begin{aligned}
 \Omega_H(M)&=\{w\in W_{G, a}:
       \Cent(w(A),{}^LH)\longrightarrow W_F\text{ is surjective}\},\\
 \mathcal W_H(M)&=W_H\backslash\Omega_H(M)/W_{M,a}.
 \end{aligned}
\]
For $w$ in this finite set, represented by
$g\in N_{\widehat G}(\widehat T_G)$, let
$P_{H,w}=M_{H,w}N_{H,w}$ be standard in $H$ with dual parabolic
and Levi $\widehat H$-conjugate to
\[
 g\,{}^LP\,g^{-1}\cap{}^LH,
 \qquad g\,{}^LM\,g^{-1}\cap{}^LH.
\]
Choose $h_w\in\widehat H$ carrying the standard dual parabolic
and Levi to these two intersections. In the identification of ${}^LH$
with its image, the inherited embedding is
\[
 \eta_{M,w}=\operatorname{Int}(g^{-1}h_w)|_{{}^LM_{H,w}}:
                   {}^LM_{H,w}\hookrightarrow{}^LM,
\]
up to $\widehat M$-conjugacy. With inherited normalizations, Xu's identity is
\[
 r_P^G\operatorname{Tran}_{H,\eta}^{G \rtimes a}(S)
  =\sum_{w\in\mathcal W_H(M)}
     \operatorname{Tran}_{M_{H,w},\eta_{M,w}}^{M\rtimes a}
                       (r_{P_{H,w}}^H S).
\]
Each double coset is counted once. Induction-transfer holds separately
for each inherited embedding, by
\cite{Xu2017-MoeglinParametrizationArthurPackets}*{
Appendix A, diagram (A.2) and the identities immediately following it}.
This termwise identity also covers the inherited data that are
nonelliptic in $M$; it is the statement in Xu's cited formula.
Thus
\[
 I_P^G r_P^G\operatorname{Tran}_{H,\eta}^{G \rtimes a}(S)
  =\operatorname{Tran}_{H,\eta}^{G \rtimes a}
       \left(\sum_{w\in\mathcal W_H(M)}
                     I_{P_{H,w}}^H r_{P_{H,w}}^H S\right).
\]
Multiply by $(-1)^{\dim A_M^a}$ and sum over $P$.
For each standard $P_H=M_HN_H$, let $a_{M_H,H,M}$ count the
double cosets $w\in\mathcal W_H(M)$ in the preceding sum for which
$P_{H,w}=P_H$. The coefficient on the right is
\[
 \sum_P(-1)^{\dim A_M^a}a_{M_H,H,M}
       =(-1)^{\dim A_{M_H}},
\]
by \cite{Xu2017-MoeglinParametrizationArthurPackets}*{Proposition A.1, equation (A.3)}.
This is exactly $\mathcal D_H$ because $H$ is semisimple. This proves
\cite{Xu2017-MoeglinParametrizationArthurPackets}*{diagram (A.1)},
whose left vertical arrow is the operator on stable distributions.
\end{proof}


\section{Preliminaries for $\rG_2$ and $\PGSO_8$}\label{sec:preliminaries-G2-PGSO8}

We now fix the root conventions and collect the facts needed for the groups $\rG_2$ and $\PGSO_8$.

\subsection{Explicit description}\label{subsec:pinned-groups}

Throughout the rest of the paper, $\rG_2$ denotes the split exceptional
group, and we put
\[
 G=\SO_8,\qquad G_{\mathrm{ad}}=\PGSO_8=G/Z(G),\qquad G_{\mathrm{sim}}=\GSO_8.
\]
All these groups are split over $F$. The notation $H$ is reserved for
an endoscopic group of $\rG_2$. 

We fix compatible pinnings and Whittaker data, with Haar measures and transfer-factor normalizations as in
Section~\ref{subsec:transfer-normalization}. 
We use the standard hyperbolic basis $e_1, \dots, e_8$ for the isotropic quadratic form on $F^8$ defining $G$ and $G_{\mathrm{sim}}$ to describe the standard maximal torus and Borel. 

Write
\[
 \begin{aligned}
 B_{\rG_2}&=T_{\rG_2}U_{\rG_2},& B_G&=T_GU_G,\\
 B_{G_{\mathrm{ad}}}&=T_{G_{\mathrm{ad}}}U_{G_{\mathrm{ad}}},& B_{G_{\mathrm{sim}}}&=T_{G_{\mathrm{sim}}}U_{G_{\mathrm{sim}}}
 \end{aligned}
\]
for the pinned Borel subgroups; these are the minimal parabolics
$P_{0,\rG_2}$, $P_{0,G}$, $P_{0,G_{\mathrm{ad}}}$, and $P_{0,G_{\mathrm{sim}}}$.
Let
\[
 j:G\hookrightarrow G_{\mathrm{sim}},\qquad
 p:G_{\mathrm{sim}} \longrightarrow G_{\mathrm{ad}},\qquad q=p\circ j:G\longrightarrow G_{\mathrm{ad}}
\]
be the inclusion and the quotients by the scalar centers. We use $q^*$
for pullback of representations of $G_{\mathrm{ad}}(F)$ to $G(F)$ and $p^*$
for inflation to $G_{\mathrm{sim}}(F)$. When explicitly stated, we use the same
symbol for a representation or distribution of $G_{\mathrm{ad}}(F)$ and its
inflation along $p$.

We identify the pinned dual groups with
\[
 \begin{aligned}
 \widehat{\rG_2}&=\rG_2(\C),& \widehat G&=\SO_8(\C),\\
 \widehat{G_{\mathrm{ad}}}&=\Spin_8(\C),& \widehat{G_{\mathrm{sim}}}&=\GSpin_8(\C).
 \end{aligned}
\]
The standard representation of $\Spin_8(\C)$ is denoted by $\Std_8$.

On the standard torus of $G_{\mathrm{sim}}$, write
\[
 t=\diag(t_1,t_2,t_3,t_4,\lambda/t_4,\lambda/t_3,\lambda/t_2,\lambda/t_1).
\]
The simple roots, which descend to $T_{G_{\mathrm{ad}}}$, are
\[
 \alpha_1(t)=t_1/t_2,\quad \alpha_2(t)=t_2/t_3,\quad
 \alpha_3(t)=t_3/t_4,\quad \alpha_4(t)=t_3t_4/\lambda.
\]
The central node is $\alpha_2$ and the outer nodes are
$\alpha_1,\alpha_3,\alpha_4$; these $D_4$ roots are distinct from
the $\rG_2$ roots $\alpha,\beta$ below.



The pinned diagram automorphisms of $G_{\mathrm{ad}}$ form $S_3$.
Fix triality $\theta$ of order three, so $A_3=\langle\theta\rangle$.
For an $S_3$-invariant irreducible $\pi'$ of $G_{\mathrm{ad}}(F)$,
$\widetilde{\pi'}$ denotes its canonical extension to
$G_{\mathrm{ad}}(F)\rtimes A_3$, as in Proposition~\ref{prop:canonical}.
Thus $\widetilde{\pi'}$ is a triality extension; the notation
$\widetilde\pi$ in the similitude discussion denotes a representation
of $G_{\mathrm{sim}}(F)$. We distinguish the twisted coset $G_{\mathrm{ad}}(F) \rtimes \theta$
from the disconnected group $G_{\mathrm{ad}}(F)\rtimes\langle\theta\rangle$.
Let 
\begin{align*}
      \widetilde{G_{\mathrm{ad}}} = \widetilde{\PGSO_8}
\end{align*}
denote the coset $G_{\mathrm{ad}}(F) \rtimes \theta$ with the twisted conjugation action of $G_{\mathrm{ad}}(F)$.

For the triality datum we put
\[
 \mathcal T=\operatorname{Tran}_{G_2}^{G_{\mathrm{ad}} \rtimes \theta, *},
 \qquad d(\theta)=|3|_F.
\]

For a parabolic $R$ of ${\rG_2}$, we suppress the superscript ${\rG_2}$
in $I_R^{\rG_2},r_R^{\rG_2},J_R^{\rG_2}$. In particular, this defines
$I_P,I_Q,J_P,J_Q$ used below.

Let $\alpha$ and $\beta$ be the short and long simple roots
of ${\rG_2}$ relative to $B_{\rG_2}$. For $\gamma\in\{\alpha,\beta\}$,
let $P_\gamma$ be generated by $B_{\rG_2}$ and the negative root
subgroup $U_{-\gamma}$. Thus $\gamma$ is the simple root
retained in its Levi. Put $P=P_\alpha$ and $Q=P_\beta$,
with Levi subgroups $M_P$ and $M_Q$.
Choose coordinates $t(t_1,t_2)$ on $T_{\rG_2}(F)$ by
\[
 \alpha(t(t_1,t_2))=t_1t_2^{-1},
 \qquad
 \beta(t(t_1,t_2))=t_2^2t_1^{-1}.
\]
Our identifications $M_P\simeq\GL_2$ and $M_Q\simeq\GL_2$
carry this element, respectively, to
\[
 \diag(t_1,t_2),\qquad
 \diag(t_2,t_1t_2^{-1}).
\]
All $\GL_2$ inducing representations on these Levis use these
identifications. In particular, the two determinant characters
on $T_{\rG_2}(F)$ are $t_1t_2$ and $t_1$, respectively, and
\[
 \delta_P|_{M_P}=\nu_2^3,\qquad
 \delta_Q|_{M_Q}=\nu_2^5.
\]
Root lengths in $\widehat{\rG_2}$ refer to the dual root system:
$\alpha^\vee$ is long and $\beta^\vee$ is short.
The names $P,Q$ always refer to the primal parabolics just fixed.

On either maximal Levi a twist $\tau\nu_2^s$ is interpreted using
the displayed $\GL_2$ identification. In particular, the notation
$\pi(a,b)$ of Section~\ref{subsec:general-representations} always
uses the intrinsic diagonal coordinates of $\GL_2$.



\subsection{Comparison with Qin's convention on $\rG_2$}\label{subsec:Qin-notation}

For comparison with \cite{Qin2026-AubertZelevinskyInvolutionG2HeckeAlgebras}, we record the
conversion between ambient $\rG_2$ torus characters and the intrinsic $\GL_2$ notation of Section~\ref{subsec:general-representations}.

Let $I_T^\gamma(a\otimes b)$ denote normalized induction to $M_{P_\gamma}(F)$ of the character $t(t_1,t_2)\mapsto a(t_1)b(t_2)$ of $T_{\rG_2}(F)$.
The torus identifications in Section~\ref{subsec:pinned-groups} give
\[
 I_T^\alpha(a\otimes b)=\pi(a,b),
 \qquad I_T^\beta(a\otimes b)=\pi(ab,a).
\]
Indeed, on $M_Q$ the intrinsic character $(a,b)$ pulls back as
$a(t_2)b(t_1/t_2)=b(t_1)(ab^{-1})(t_2)$.
Thus, intrinsic labels $(a,b)$ on $M_Q$ correspond to ambient labels
$b\otimes ab^{-1}$. 
The determinant characters computed above imply
\[
 \begin{aligned}
 I_T^\alpha(a\otimes b)\nu_2^s
    &=I_T^\alpha(\nu^sa\otimes\nu^sb),\\
 I_T^\beta(a\otimes b)\nu_2^s
    &=I_T^\beta(\nu^sa\otimes b).
 \end{aligned}
\]

In the induction notation $I_\gamma(s,\tau)$ and quotient notation
$J_\gamma(s,\tau)$ used for these computations in
\cite{Qin2026-AubertZelevinskyInvolutionG2HeckeAlgebras}, once $\tau$
has been expressed in intrinsic $\GL_2$ coordinates, the dictionary is
\[
 \begin{array}{c|cc}
 &\gamma=\alpha&\gamma=\beta\\ \hline
 I_\gamma(s,\tau)&I_P(\tau\nu_2^s)&I_Q(\tau\nu_2^s)\\
 J_\gamma(s,\tau)&J_P(\tau\nu_2^s)&J_Q(\tau\nu_2^s).
 \end{array}
\]
The second row applies to Langlands standard modules, and the
Steinberg twist denoted $\delta(\chi)$ there is our $\St_2(\chi)$.

In particular, if $\chi^3=1$ and $\chi\ne1$, then ambient labels
$\chi^{-1}\otimes\chi^{-1}$ on $M_Q$ give
\[
 I_T^\beta(\chi^{-1}\otimes\chi^{-1})
     =\pi(\chi^{-2},\chi^{-1})=\pi(\chi,\chi^{-1}).
\]
Throughout this paper the corresponding cubic quotient is therefore
written $J_Q(\pi(\chi,\chi^{-1})\nu_2)$.
The symbol $\pi(a,b)$ always retains its intrinsic $\GL_2$ meaning.

\section{Representation theory of $\rG_2(F)$}

\subsection{Elliptic endoscopic groups of $\rG_2$}\label{subsec:elliptic-endoscopic-groups}

The elliptic endoscopic groups of $\rG_2$ are  given by $\rG_2$, $\SO_4$ and $\PGL_3$.
Let $\sn$ be the spinor norm map 
\begin{align}
     \sn \colon \SO_4(F) \to F^\times / F^{\times 2}.
\end{align}
We also have the determinant map
\begin{align}
      \det \colon \PGL_3(F) \to F^\times / F^{\times 3}.
\end{align}

Let $P_l$ (resp. $P_s$) be the parabolic subgroup of $\SO_4$ corresponding to the short root (resp. the long root) dual parabolic subgroup in $\widehat{\rG_2}(\C)$.

\subsection{\texorpdfstring{Arthur parameters for $\rG_2$ and their standard lifts}{Arthur parameters for G2 and their standard lifts}}\label{subsec:G2-parameters}

We apply Definition~\ref{def:parameters} to $\rG_2$.
Since $Z(\widehat{\rG_2})=1$, the component group used for ${\rG_2}$ is
\[
 \cS_\psi=\pi_0\Cent_{\widehat{\rG_2}}(\operatorname{Im}\psi).
\]
For a nontrivial Arthur $\SL_2$, its conjugacy class and centralizer
are as follows:
\begin{center}
\begin{tabular}{c|c}
Arthur $\SL_2$ class & centralizer in $\widehat{\rG_2}$\\ \hline
short-root & long-root $\SL_2$\\
long-root & short-root $\SL_2$\\
subregular & $S_3$\\
principal & $1$
\end{tabular}
\end{center}
The table is the list of nonzero nilpotent-orbit types and their
$\SL_2$-centralizers in
\cite{GanSavin2023-HoweDualityDichotomyExceptionalThetaCorrespondences}*{Section 2.7}.
The trivial Arthur factor gives the tempered case. In the subregular
case the connected Deligne factor maps trivially to the finite
centralizer, so $\psi$ is co-tempered. The possible subgroups
$A=\psi(W_F)\subset S_3$, up to conjugacy, are
$1$, $\mu_2$, $\mu_3$, and $S_3$.
For principal Arthur image the centralizer is trivial, so the Weil
and Deligne factors are trivial. 
For long-root and short-root parameters, the Deligne factor maps to the opposite-root $\SL_2$, either trivially or surjectively. 
The former case is co-tempered; in the latter the Weil image lies in the center $\mu_2$ of this $\SL_2$. 

\begin{definition}
      In the remaining cases, the image of the Deligne and Arthur $\SL_2$ appears in the standard representation as $S_2^D\boxtimes S_2^A$ or $S_2^A\boxtimes S_2^D$, respectively, where $S_n$ is the $n$-dimensional irreducible representation of $\SL_2(\C)$.
      We call these parameters mixed parameters.
\end{definition}


Fix the standard embedding $\iota = \iota^{\PGSO_8}_{\rG_2} : \widehat{\rG_2}\hookrightarrow\widehat{G_{\mathrm{ad}}}=\Spin_8(\C)$.
Put $\psi'=\iota\circ\psi$. The dual maps of $q$ and $p$ are
$\widehat q:\widehat{G_{\mathrm{ad}}}\to\widehat G$ and
$\widehat p:\widehat{G_{\mathrm{ad}}}\hookrightarrow\widehat{G_{\mathrm{sim}}}$.
We write $\psi_0=\widehat q\circ\psi'=\Std_8\circ\psi'$ for
the orthogonal parameter and $\widetilde\psi=\widehat p\circ\psi'$
for its similitude lift. For the two
commuting root $\SL_2$ factors, the standard representation restricts as
\[
 \Std_8=(S_{2,s}\boxtimes S_{2,l})\oplus S_{3,s}\oplus\mathbf1.
\]
Here $s$ and $l$ specify the short and long factors in $\widehat{\rG_2}$.

\begin{corollary}\label{cor:mixed-parameters}
For a quadratic character $\chi$, including $\chi=1$, the long-root
Arthur parameter with Deligne parameter $\chi\boxtimes S_2$ in the
short-root centralizer has standard lift
\begin{align*}
 \Std_8\circ\psi'_D&=\chi\boxtimes S_2^D\boxtimes S_2^A
                       \oplus S_3^D\oplus\mathbf1,\\
 \Std_8\circ\phi_{\psi'_D}&=
  \chi\nu^{1/2}\boxtimes S_2^D\oplus\chi\nu^{-1/2}\boxtimes S_2^D
                       \oplus S_3^D\oplus\mathbf1.
\end{align*}
For the swapped, short-root parameter the corresponding expressions are
\begin{align*}
 \Std_8\circ\psi'_A&=\chi\boxtimes S_2^D\boxtimes S_2^A
                       \oplus S_3^A\oplus\mathbf1,\\
 \Std_8\circ\phi_{\psi'_A}&=
  \chi\nu^{1/2}\boxtimes S_2^D\oplus\chi\nu^{-1/2}\boxtimes S_2^D
                       \oplus\nu\oplus\mathbf1\oplus\nu^{-1}\oplus\mathbf1.
\end{align*}
\end{corollary}
\begin{proof}
Use the restriction of $\Std_8$ above and $\chi^2=1$, and then apply
the block formula~\eqref{eq:parameter-block}. Swapping the factors interchanges the
$S_3^D$ and $S_3^A$ summands.
\end{proof}

Subregular parameters factor through
\begin{align*}
      L_F \times \SL_2^A \to W_F \times \SL_2^A \xrightarrow{\sigma \times S_3^A} S_3 \times \SO_3 \hookrightarrow \rG_2.
\end{align*}
The action of $S_3 \times \SO_3$ on the standard representation of $\rG_2$ is given by
\begin{align*}
      \rho \boxtimes S_3^A \oplus \epsilon \boxtimes S_1^A,
\end{align*}
where $\rho$ is the standard 2-dimensional representation of $S_3$ and $\epsilon$ is the sign character of $S_3$.

\begin{lemma}\label{lemma:subregular-standard}
      The lift $\psi'_{\sigma} = \iota \circ \psi_{\sigma}$ of a subregular A-parameter
      \begin{align*}
            \psi_{\sigma} \colon L_F \times \SL_2^A \to W_F \times \SL_2^A \xrightarrow{\sigma \times S_3^A} S_3 \times \SO_3 \hookrightarrow \rG_2.
      \end{align*}
      to $\PGSO_8$ has the standard parameter
      \begin{align*}
           \Std_8 \circ \psi'_{\sigma} = \rho \circ \sigma \boxtimes S_3^A \oplus \det \rho \circ \sigma \oplus \mathbf{1}.
      \end{align*}
      The swapped parameter is a tempered $L$-parameter
      \begin{align*}
            \rho \circ \sigma \boxtimes S_3^D \oplus \det \rho \circ \sigma \oplus \mathbf{1}.
      \end{align*}
\end{lemma}


\subsection{A rank-one Aubert calculation}

\begin{lemma}\label{lem:rank-one}
Let $R$ be $P$ or $Q$, and let $\tau$ be unitary supercuspidal on its
$\GL_2$ Levi. Suppose the standard module $I_R(\tau\nu_2^s)$, $s>0$,
has length two, with a square-integrable subrepresentation $\delta$ and
Langlands quotient $J$. Then
\[
 [J]=[I_R(\tau\nu_2^s)]-[\delta],\qquad
 \mathcal D_{\rG_2}[\delta]=-[J],\qquad\widehat\delta=J.
\]
\end{lemma}
\begin{proof}
The geometric lemma gives the two cuspidal exponents $\tau\nu_2^s$
and its Weyl conjugate in $r_R I_R(\tau\nu_2^s)$, each with
multiplicity one. Frobenius reciprocity supplies the first exponent in
$r_R\delta$; the square-integrability criterion excludes the opposite
one. Thus $r_R\delta=\tau\nu_2^s$. Its Jacquet modules for the torus
and for the nonconjugate maximal Levi vanish, since its supercuspidal
support is in $M_R$. Formula~\eqref{eq:Aubert-operator} now gives
$\mathcal D_{\rG_2}[\delta]=[\delta]-[I_R(\tau\nu_2^s)]=-[J]$.
The support Levi has semisimple rank one, so
$\epsilon_{\rG_2}(\delta)=-1$ and Proposition~\ref{prop:Aubert} gives
the last assertion. The exact sequences used below are supplied by
\cite{GanSavin2023-HoweDualityDichotomyExceptionalThetaCorrespondences}*{
Propositions 3.1(i) and 3.2(i)}.
\end{proof}



\section{Representation theory of $\widetilde{\PGSO_8}(F)$}


\subsection{Similitude packets and the distinguished tempered lift}\label{subsec:similitude-preliminaries}

Let $\Sigma_0=\langle\theta_0\rangle$, where $\theta_0$ is the standard pinned outer involution of $G$, extended to $G_{\mathrm{sim}}$;
we use the corresponding involutions on the dual groups.
Bars on $\Irr$ and packet symbols denote passage to $\Sigma_0$-orbits,
and brackets denote these orbits on representations or parameter
conjugacy classes, using the dual involution for parameters.

We use the infinitesimal character of
Section~\ref{subsec:parameters-and-infinitesimal-characters}, in particular
$\lambda_{\widetilde\psi}(w)=\widetilde\psi(w,d_w,d_w)$.
Fix the local Langlands correspondence modulo $\Sigma_0$ with the
choices of \cite{Xu2024-ArthurPacketsQuasisplitGSp2nGO2n}*{Section 5}.
For $[\widetilde\pi]\mapsto[\widetilde\varphi_{\widetilde\pi}]$,
put $[\lambda_{\widetilde\pi}]
=[\lambda_{\widetilde\varphi_{\widetilde\pi}}]$ and define
\begin{equation}\label{eq:Irr-infinitesimal-character}
 \overline{\Irr}(G_{\mathrm{sim}}(F))_{\widetilde\lambda}
 =\left\{[\widetilde\pi]\in\Irr(G_{\mathrm{sim}}(F))/\Sigma_0:
       [\lambda_{\widetilde\pi}]=[\widetilde\lambda]\right\}.
\end{equation}
Thus the subscript fixes the infinitesimal character in the sense of
\cite{Xu2024-ArthurPacketsQuasisplitGSp2nGO2n}*{Section 3,
pp.~10353--10356}.

For $\widetilde\xi=\widetilde\lambda$ or $\widetilde\psi$, write
\[
 X^{\Sigma_0}(\widetilde\xi)=
 \left\{\omega\in\Hom(G_{\mathrm{sim}}(F)/G(F),\C^\times):
       [\widetilde\xi\otimes\omega]=[\widetilde\xi]\right\}.
\]
Here $\omega$ is a smooth character, identified with a character
of $W_F$ by local class field theory, and the twist uses the central $\C^\times$ in
$\widehat{G_{\mathrm{sim}}}\to\widehat G$.

\begin{proposition}\label{prop:Xu-packets}
Fix a central character of $G_{\mathrm{sim}}(F)$ extending that of a
tempered packet of $G(F)$. The refined
similitude packets of \cite{Xu2018-LPacketsQuasisplitGSp2nGO2n}*{Theorem 4.6}
have stable sums of characters and satisfy the character relations of
that theorem. Their choice is unique only up to the character group
\[
 X=\Hom\bigl(G_{\mathrm{sim}}(F)/(Z_{G_{\mathrm{sim}}}(F)G(F)),\C^\times\bigr).
\]
For Arthur parameters, the corresponding existence, stability, and
character identities are supplied by
\cite{Xu2024-ArthurPacketsQuasisplitGSp2nGO2n}*{Theorem 6.4}, with its
central-character, outer-conjugacy, and normalization conventions.
\end{proposition}
\begin{proof}
The tempered assertion follows from
\cite{Xu2018-LPacketsQuasisplitGSp2nGO2n}*{Theorem 4.6(1)--(3)},
and the twisting ambiguity is described in
\cite{Xu2018-LPacketsQuasisplitGSp2nGO2n}*{Remark 4.7}.
The Arthur assertion is
\cite{Xu2024-ArthurPacketsQuasisplitGSp2nGO2n}*{Theorem 6.4(1)--(3)}.
\end{proof}

For $G_{\mathrm{ad}}$, we impose the trivial scalar central character on
$G_{\mathrm{sim}}(F)$. If $\pi'$ is a representation of $G_{\mathrm{ad}}(F)$, the
similitude parameter of $p^*\pi'$ therefore factors through
$\Spin_8(\C)\subset\GSpin_8(\C)$. Write $\phi_{\pi'}$ for a representative
of this parameter modulo the standard outer involution.
The linear parameter $\Std_8\circ\phi_{\pi'}$ is independent of this
choice up to isomorphism. In the following criterion,
$\gamma(s,r,\psi_F)$ is the Artin gamma factor of the linear
Weil--Deligne representation $r$; the parameter of $s_1\pi'$ is
obtained separately in the same way. We use the following
distinguished tempered packet.

\begin{theorem}\label{thm:distinguished}
Let $s_1$ be the pinned involution interchanging the standard and one
half-spin representation of $\Spin_8$. A tempered irreducible $\pi'$ of
$G_{\mathrm{ad}}(F)$ satisfying
\[
 \gamma(0,\Std_8\circ\phi_{\pi'},\psi_F)=
 \gamma(0,\Std_8\circ\phi_{s_1\pi'},\psi_F)=0
\]
is the composite theta lift of a unique tempered representation of ${\rG_2}(F)$.
The map taking $\pi'$ to the parameter of this representation is surjective
onto the tempered parameters of ${\rG_2}$. Its fibers are distinguished Xu
packets, denoted $\Pi^X_{\iota\circ\varphi}$. Every member is invariant
under all pinned diagram automorphisms $S_3$ of $G_{\mathrm{ad}}$.
\end{theorem}
\begin{proof}
The characterization and the assertion about fibers are
\cite{TakanashiThesisLocal}*{Proposition 3.58 and Definition 3.59,
pp.~25--26}. The $S_3$-invariance follows from
\cite{TakanashiThesisLocal}*{Proposition 6.6}. The underlying
tempered correspondence for ${\rG_2}$ is
\cite{GanSavin2023-LocalLanglandsConjectureG2}; the similitude packets
are those of Proposition~\ref{prop:Xu-packets}.
\end{proof}


\subsection{Restriction and compatibility with Arthur packets}

We use the maps $j,p,q$ of Section~\ref{subsec:pinned-groups} and the positive
Aubert functor $D_A$ of Section~\ref{subsec:aubert-involution-and-endoscopic-transfer}.
For $G_{\mathrm{ad}}(F)\rtimes S_3$, the same symbol denotes the positive
disconnected Aubert functor of
\cite{ChengKaletha2026-AubertBernsteinInvolutionsDisconnectedGroups}*{Definition 3.1.9}.

\begin{lemma}\label{lem:Aubert-restriction}
Restriction from $G_{\mathrm{sim}}(F)$ to $G(F)$, pullback along $q$,
and restriction from $G_{\mathrm{ad}}(F)\rtimes S_3$ to $G_{\mathrm{ad}}(F)$ commute with $D_A$.
\end{lemma}
\begin{proof}
Corresponding standard parabolics $\widetilde R$ and $R$ of
$G_{\mathrm{sim}}$ and $G$ have the same unipotent radical and modulus
character, and
$G_{\mathrm{sim}}(F)=\widetilde R(F)G(F)$. Thus restriction commutes with
normalized induction, Jacquet modules, and the maps in the Aubert
complex. The relative semisimple ranks agree, so taking its cohomology
gives the assertion for $D_A$. Inflation from $G_{\mathrm{ad}}$ to $G_{\mathrm{sim}}$
also identifies these complexes, proving the assertion for $q$.
The last assertion is
\cite{ChengKaletha2026-AubertBernsteinInvolutionsDisconnectedGroups}*{
Lemma 3.1.10}.
\end{proof}

\begin{proposition}\label{prop:standard-Aubert}
Let $\psi_0=\Std_8\circ\psi'$. If
$\operatorname{JH}(q^*\pi')\subset\Pi_{\psi_0}(G)$, then
\[
 \operatorname{JH}(q^*\widehat{\pi'})
       \subset\Pi_{\widehat{\psi_0}}(G),\qquad
 \widehat{\psi_0}=\Std_8\circ\widehat{\psi'}.
\]
Packets of $G$ are taken modulo its standard outer involution.
\end{proposition}
\begin{proof}
For classical Arthur packets,
$\widehat{\Pi_{\psi_0}}=\Pi_{\widehat{\psi_0}}$ by
\cite{AtobeGanEtAl2025-LocalIntertwiningRelationsCoTemperedPackets}*{
Lemma 4.4.4(1)}, whose character-relation hypotheses follow from
Arthur's local classification, as recalled in
\cite{Xu2017-MoeglinParametrizationArthurPackets}*{Theorems 4.1--4.2}.
In the even orthogonal case the cited lemma is formulated for the
full orthogonal group $\mathrm O_8(F)$. Restricting its packet members
to $G(F)=\SO_8(F)$ recovers the classical packets modulo $\Sigma_0$;
see \cite{AtobeGanEtAl2025-LocalIntertwiningRelationsCoTemperedPackets}*{
Remark 1.6.2(1)}.
Positive Aubert duality commutes with restriction to the identity
component by
\cite{ChengKaletha2026-AubertBernsteinInvolutionsDisconnectedGroups}*{
Lemma 3.1.10}, so the same packet identity holds in these
outer-involution orbits. Apply Lemma~\ref{lem:Aubert-restriction}
to $q^*\pi'$. The parameter equality follows
by interchanging the two $\SL_2$ factors.
\end{proof}

\begin{proposition}\label{prop:Xu-Aubert}
Let $\widetilde\psi$ be the $G_{\mathrm{sim}}$ parameter induced by $\psi'$,
and put $\widetilde\lambda=\lambda_{\widetilde\psi}$. In Xu's notation,
suppose
\begin{equation}\label{eq:Xu-criterion}
 X^{\Sigma_0}(\widetilde\lambda)
       =X^{\Sigma_0}(\widetilde\psi).
\end{equation}
Then the positive Aubert duals of the Xu packet for $\psi'$ form the
Xu packet for $\widehat{\psi'}$, with the same infinitesimal and central
characters. 
Condition~\eqref{eq:Xu-criterion} holds if $\psi'$ is
tempered or co-tempered.

Furthermore, in this case the restriction map induces a bijection
\begin{align*}
      \overline{\Pi}_{\widetilde{\psi}}(G_{\mathrm{sim}}) &\longrightarrow \overline{\Pi}_{\psi_0}(G)/G_{\mathrm{sim}}(F).
\end{align*}
\end{proposition}

\begin{proof}
By \cite{Xu2024-ArthurPacketsQuasisplitGSp2nGO2n}*{Theorem 6.5
and the preceding construction}, the packet is
\[
 \overline\Pi_{\widetilde\psi}
 =\left\{[\widetilde\pi]\in
       \overline{\operatorname{Irr}}(G_{\mathrm{sim}}(F))_{\widetilde\lambda}:
       \operatorname{JH}(\widetilde\pi|_{G(F)})
                         \subset\overline\Pi_{\psi_0}\right\}.
\]
The notation is as in Section~\ref{subsec:similitude-preliminaries}.
Aubert duality preserves cuspidal support, hence the infinitesimal
character by
\cite{Xu2024-ArthurPacketsQuasisplitGSp2nGO2n}*{Proposition 5.1 and Corollary 5.5};
it also preserves
the central character. Moreover,
\[
 \lambda_{\widehat{\widetilde\psi}}=\widetilde\lambda,
 \qquad
 X^{\Sigma_0}(\widehat{\widetilde\psi})
       =X^{\Sigma_0}(\widetilde\psi),
\]
since swapping the two $\SL_2$ factors changes neither the diagonal
parameter nor the image. Proposition~\ref{prop:standard-Aubert} gives
one inclusion between the packets; involutivity gives the reverse.
For a tempered or co-tempered parameter the diagonal parameter
$\widetilde\psi_d(w,x)=\widetilde\psi(w,x,x)$ has the same image as
$\widetilde\psi$ and is the open-orbit parameter for
$\widetilde\lambda$. By
\cite{Xu2024-ArthurPacketsQuasisplitGSp2nGO2n}*{Section 3, p.~10356},
we have
$X^{\Sigma_0}(\widetilde\psi_d)=X^{\Sigma_0}(\widetilde\lambda)$,
proving \eqref{eq:Xu-criterion}.
\end{proof}




\section{Tempered twisted endoscopy for $\widetilde{\PGSO_8}$ and Aubert duality}\label{sec:twisted-endoscopy}

\subsection{Tempered character identities and stability on proper Levi subgroups}

For a tempered parameter $\varphi$ of $\rG_2$, write $\Pi_\varphi$
for its Gan--Savin packet and $\rho_{\varphi,\pi}\in\Irr(\cS_\varphi)$ for the enhancement in the fixed Whittaker normalization constructed in \cite{TakanashiEndoscopyLocal}*{Theorem 31.1}.
Set
\[
 m_\pi=\dim\rho_{\varphi,\pi},\qquad
 S\Theta_\varphi=\sum_{\pi\in\Pi_\varphi}m_\pi\Theta_\pi.
\]
This is the stable tempered distribution used below.

\begin{theorem}\label{thm:tempered-input}
For each tempered parameter $\varphi$ considered below, one has
\begin{equation}\label{eq:tempered-triality}
 \mathcal T(S\Theta_\varphi)=
 \sum_{\pi'\in\Pi^X_{\iota\circ\varphi}}\Theta_{\widetilde{\pi'}}.
\end{equation}
More generally, we also have the twisted endoscopic character identities
\begin{align*}
      \sum_{\pi'\in\Pi^X_{\iota\circ\varphi}} \langle s, \pi' \rangle \Theta_{\widetilde{\pi'}}(f^{\widetilde{\PGSO_8}})
      =
      S\Theta^{H_s}_{\varphi^s}(\Trans_{H_s}^{\widetilde{\PGSO_8}} f^{\widetilde{\PGSO_8}}),
\end{align*}
for all $s\in\cS_{\iota_{\rG_2}^{\PGSO_8}\circ\varphi}$, where $\varphi^s$ is the
parameter of the endoscopic group $H_s$ corresponding to $s$.

The tempered endoscopic identities for $\rG_2$ used in
Sections~\ref{sec:subregular}--\ref{sec:mixed-short} hold with the same
normalization.
\end{theorem}

\begin{proof}
This is the non-archimedean case of
\cite{TakanashiEndoscopyLocal}*{Theorem 31.1},
applied to the distinguished packets of Theorem~\ref{thm:distinguished}.
\end{proof}

\begin{remark}
      In \cite{TakanashiEndoscopyLocal}, we assumed the validity of the (twisted) weighted fundamental lemma for $\widetilde{\PGSO_8}$ and $\rG_2$.
      A proof of the general weighted fundamental lemma is announced in \cite{HalleckDube2026-WeightedFundamentalLemmaNonSplitGroups}.
\end{remark}

The proper Levis on the $D_4$ side have the following additional
stability property.

\begin{lemma}\label{lem:proper-Levi-stable}
Every finite-length character of a proper Levi subgroup of $G_{\mathrm{sim}}$
or $G_{\mathrm{ad}}$ is stable. Normalized induction from a proper Levi
therefore preserves stability without any condition on the inducing
finite-length representation.
\end{lemma}
\begin{proof}
Such a Levi is a quotient $L/Z$ of a product of general linear groups
by a central split torus. Indeed, up to the outer involution, proper
Levis of $G_{\mathrm{sim}}$ are products of general linear groups and either
a split torus or $\GSO_{2m}$ with $m\leq3$, and
\[
 \GSO_4\simeq(\GL_2\times\GL_2)/\{(zI_2,z^{-1}I_2)\},\qquad
 \GSO_6\simeq(\GL_4\times\GL_1)/\{(zI_4,z^{-2})\}.
\]
Passing to $G_{\mathrm{ad}}$ quotients by another central split torus.
Lemma~\ref{lem:central-quotient-stability} therefore gives stability
of every finite-length character of such a Levi. Apply
Lemma~\ref{lem:induction-transfer} for the last assertion.
\end{proof}

We fix the inherited Levi embeddings for the later calculations.
For $R=P,Q$ with $\GL_2$ Levi $M_R$, let $M_{R'}$ be the
$\theta$-stable Levi of $G_{\mathrm{ad}}$ whose dual group is
\[
 \widehat M_{R'}=\Cent(\iota(Z(\widehat M_R)^\circ),\Spin_8(\C)).
\]
Take the standard parabolic $R'=M_{R'}N_{R'}$ with this Levi.
Restriction of $\iota$ gives the inherited Levi embedding; denote
its transfer by $\mathcal T_{M_R}^{M_{R'},\theta}$. The derived types
of $M_{P'},M_{Q'}$ are $A_1,A_1^3$, respectively. Thus
Lemma~\ref{lem:induction-transfer} gives, for stable virtual characters $S$,
\[
      \mathcal T(I_R S)=I_{R'}^{G_{\mathrm{ad}}}
      \mathcal T_{M_R}^{M_{R'} \rtimes \theta}(S).
\]

The parabolic subgroup $Q' = P_{\alpha_1, \alpha_3, \alpha_4}$ of $\PGSO_8$ has Levi subgroup $M_{Q'} = M_{\alpha_1, \alpha_3, \alpha_4} = (\GL_2^3 \times \GL_1)/ \GL_1^3$, where the embedding is given by $(z_1, z_2, z_3) \mapsto (z_1, z_2, z_3, z_1 z_2 z_3)$.
The action of the triality group $S_3$ on $M_{Q'}$ is given by permuting the three $\GL_2$-factors.
For any representation $\sigma$ of $M_Q \cong \GL_2$, we define a  representation of $M_{Q'}$ by
\begin{align*}
      \sigma \boxtimes \sigma \boxtimes \sigma \boxtimes \omega_{\sigma}^{-1}
\end{align*}
and denote it by $(\sigma)_{M_{Q'}}$. 
Then, the canonical extension to $M_{Q'} \rtimes \theta$ is equal to $\Trans_{M_{Q}}^{M_{Q'} \rtimes \theta}(\sigma)$.
We have $M_{P'} \cong \GL_2 \times \GL_1^3/\GL_1$ where the embedding is given by $z \mapsto (zI_2, z, z, z)$.
We use similar notation $(\sigma)_{M_{P'}}$ for the representation of $M_{P'}$ defined by
\begin{align*}
      \sigma \omega_{\sigma} \boxtimes \omega_{\sigma}^{-1} \boxtimes \omega_{\sigma}^{-1} \boxtimes \omega_{\sigma}^{-1}.
\end{align*}
For the standard Borel subgroup $B_{\PGSO_8}$ of $\PGSO_8$, a triality-stable representation of the maximal torus $T_{\PGSO_8}$ corresponds to a unique representation of $T_{\rG_2}$ via 
\begin{align*}
      (T_{\PGSO_8})_{\theta} \cong T_{\rG_2}.
\end{align*}
For a character $\chi$ of $T_{\rG_2}$, we denote the corresponding character of $T_{\PGSO_8}$ by $(\chi)_{T_{\PGSO_8}}$.

From the Langlands classification, we obtain the following result.
\begin{lemma}\label{lem:triality-invariant-non-tempered}
      Let $(R', \sigma')$ be a pair consisting of a proper standard parabolic subgroup $R' = M_{R'} N_{R'}$ of $\PGSO_8$ and an essentially tempered representation $\sigma'$ of $M_{R'}$ such that the central character of $\sigma'$ is positive with respect to $R'$. Then, the corresponding Langlands quotient $J_{R'}(\sigma')$ of $\PGSO_8$ is triality-invariant if and only if $R' \in \{B_{\PGSO_8}, P', Q' \}$ and $\sigma'$ is triality-invariant.

      In this case, there exists a unique pair of a standard parabolic subgroup $R$ of $\rG_2$ and an essentially tempered representation $\sigma$ of $M_R$ such that  $R'$ corresponds to $R$ and $\sigma' = (\sigma)_{M_{R'}}$ corresponds to $\sigma$.
\end{lemma}



\subsection{Triality signs and Aubert transfer}

Specializing \eqref{eq:twisted-Aubert-operator} gives
\[
 \mathcal D_{G_{\mathrm{ad}}}^\theta=
  \sum_{\substack{R=MN\supset B_{G_{\mathrm{ad}}}\\\theta(R)=R}}
    (-1)^{\dim A_M^\theta}I_R^{G_{\mathrm{ad}}}r_R^{G_{\mathrm{ad}}}.
\]
The following proposition identifies its support sign on canonical
triality extensions.

\begin{proposition}\label{prop:twisted-sign}
Let $\pi'$ be $S_3$-invariant and irreducible. Then
\begin{equation}\label{eq:twisted-sign}
 \mathcal D_{G_{\mathrm{ad}}}^\theta\Theta_{\widetilde{\pi'}}
       =\epsilon_{G_{\mathrm{ad}}}(\pi')\Theta_{\widetilde{\widehat{\pi'}}}
       \quad\text{on }G_{\mathrm{ad}}(F)\theta.
\end{equation}
\end{proposition}
\begin{proof}
The Aubert dual commutes with restriction to $G_{\mathrm{ad}}(F)$ and preserves
irreducibility, up to sign, by
\cite{ChengKaletha2026-AubertBernsteinInvolutionsDisconnectedGroups}*{
Lemmas 3.1.10--3.1.11}. Apply it to an $S_3$-extension of $\pi'$.
The resulting representation is an $S_3$-extension of $\widehat{\pi'}$,
so Proposition~\ref{prop:canonical} identifies its restriction to $A_3$
with the canonical extension. The character formula and sign comparison
in \cite{ChengKaletha2026-AubertBernsteinInvolutionsDisconnectedGroups}*{
Proposition 3.2.1 and Corollary 3.2.4} give the scalar $\epsilon_{G_{\mathrm{ad}}}(\pi')$.
Note that the orientation twist in
\cite{ChengKaletha2026-AubertBernsteinInvolutionsDisconnectedGroups}*{Corollary 3.3.2}
is trivial on $\theta$.
\end{proof}

\begin{theorem}\label{thm:signed-transfer}
Suppose a stable virtual character $S=\sum_i a_i\Theta_{\pi_i}$ of ${\rG_2}$
satisfies
$\mathcal T(S)=\sum_j b_j\Theta_{\widetilde{\pi'_j}}$, with the
$\pi'_j$ irreducible and $S_3$-invariant. 

Then, we have
\begin{equation}\label{eq:signed-transfer}
 \mathcal T\left(\sum_i a_i\epsilon_{\rG_2}(\pi_i)
                         \Theta_{\widehat\pi_i}\right)
  =\sum_j b_j \epsilon_{G_{\mathrm{ad}}}(\pi'_j)
       \Theta_{\widetilde{\widehat{\pi'_j}}}.
\end{equation}
The distribution inside $\mathcal T$ is stable. The ordinary analogue for $\rG_2$ follows from the compatibility formula in Section~\ref{subsec:aubert-involution-and-endoscopic-transfer}.
\end{theorem}

\begin{proof}
Apply $\mathcal D_{G_{\mathrm{ad}}}^\theta$ to the assumed equality, move it across
$\mathcal T$ by Proposition~\ref{prop:Xu-transfer}, and use
Propositions~\ref{prop:Aubert} and \ref{prop:twisted-sign} term by term.
Stability is the left vertical assertion of
\cite{Xu2017-MoeglinParametrizationArthurPackets}*{diagram (A.1)}.
For ordinary transfer the same argument uses Hiraga's formula.
\end{proof}

\begin{corollary}\label{cor:cotempered}

For a co-tempered parameter $\psi$, put $\varphi=\widehat\psi$ and
define
\[
 S\Theta_\psi^{\mathrm{Aub}}=\mathcal D_{\rG_2} S\Theta_\varphi,
 \qquad T_\psi^{\mathrm{Aub}}=
       \mathcal D_{G_{\mathrm{ad}}}^\theta\sum_{\pi'\in\Pi^X_{\iota\circ\varphi}}
                                      \Theta_{\widetilde{\pi'}}.
\]
Then $S\Theta_\psi^{\mathrm{Aub}}$ is stable and
$\mathcal T(S\Theta_\psi^{\mathrm{Aub}})=T_\psi^{\mathrm{Aub}}$.
Its irreducible terms are the Aubert duals of those of $S\Theta_\varphi$,
with the coefficients prescribed in \eqref{eq:signed-transfer}.
On $G_{\mathrm{ad}}$ these duals form the distinguished Xu packet
\[
 \Pi^X_{\iota\circ\psi}
       =\{\widehat{\pi'}:\pi'\in\Pi^X_{\iota\circ\varphi}\}.
\]
\end{corollary}

\begin{proof}
Apply Theorem~\ref{thm:signed-transfer} to
\eqref{eq:tempered-triality} and Proposition~\ref{prop:Xu-Aubert}.
Lemma~\ref{lem:Aubert-restriction} preserves the $S_3$-extensions,
whose restrictions to $A_3$ are canonical by
Proposition~\ref{prop:canonical}.
\end{proof}

The support signs are
\[
\begin{array}{c|c|c}
 \text{group}&\text{semisimple type of the supercuspidal support Levi}&\text{support sign}\\\hline
 {\rG_2}&\text{torus}&+1\\
 {\rG_2}&A_1&-1\\
 {\rG_2}&\rG_2&+1\\
 G_{\mathrm{ad}}&\text{torus, }A_1^2\text{, or }A_2&+1\\
 G_{\mathrm{ad}}&A_1\text{, }A_3\text{, or }A_1^3&-1\\
 G_{\mathrm{ad}}&D_4&+1
\end{array}
\]


\section{The principal parameter}\label{sec:principal-parameter}

\begin{proposition}\label{prop:principal-ecr}
For the principal Arthur parameter, the Aubert-dual set is
$\{\mathbf1_{\rG_2}\}$, where $\mathbf1_{\rG_2}=J_{B_{\rG_2}}^{\rG_2}(\delta_{B_{\rG_2}}^{1/2})$, and
\[
 \mathcal T(\Theta_{\mathbf1_{\rG_2}})=\Theta_{\mathbf1_{\PGSO_8}},
\]
where the trivial representation on the right has the trivial triality
action.
\end{proposition}
\begin{proof}
The swapped parameter has packet $\{\St_{\rG_2}\}$. Its lift is principal in
$\Spin_8$: its standard representation is $S_7\oplus\mathbf1$.
The distinguished tempered lift is $\St_{G_{\mathrm{ad}}}$, and
Theorem~\ref{thm:tempered-input} gives
$\mathcal T(\Theta_{\St_{\rG_2}})=\Theta_{\widetilde{\St_{G_{\mathrm{ad}}}}}$.
The positive Aubert dual of the Steinberg representation is the trivial
representation. Both Steinberg representations have toral supercuspidal
support, and the ranks are two and four, so both signs are $+1$.
Theorem~\ref{thm:signed-transfer} proves the assertion.
\end{proof}

\section{Subregular parameters}\label{sec:subregular}


\subsection{The tempered subregular packet}\label{sec:tempered-subregular}

Let $\psi$ be a subregular Arthur parameter of ${\rG_2}$.
Then, the swapped parameter $\phi = \widehat\psi$ is tempered.
We use the description in Section \ref{subsec:G2-parameters} and write $\phi = \phi_\sigma$ for a homomorphism
\begin{align*}
      \sigma \colon W_F \to S_3.
\end{align*}
Its image is either $1$, $\mu_2$, $\mu_3$, or $S_3$.
In the second and third cases, let $\chi$ be the corresponding nontrivial quadratic or cubic character of $W_F$.

\begin{theorem}\label{thm:tempered-subregular-elements}
      We have the following tempered subregular $L$-packets of ${\rG_2}$.
      \begin{enumerate}
            \item 
            If $\sigma(W_F) = \{1\}$, then the $L$-packet $\Pi_{\phi}$ consists of three representations $\pi_{\gen}[1], \pi_{\deg}[1], \pi_{\mathrm{sc}}[1]$ corresponding to the three irreducible representations $\mathbf{1}, \rho, \epsilon$ of $S_3$.
            The representation $\pi_{\gen}[1]$ is generic, $\pi_{\deg}[1]$ is non-generic, and $\pi_{\mathrm{sc}}[1]$ is supercuspidal.
            \item 
            If $\sigma(W_F) = \mu_2$, then the $L$-packet $\Pi_{\phi}$ consists of two representations $\pi_{\gen}[\chi], \pi_{\mathrm{sc}}[\chi]$ corresponding to the two irreducible representations $\mathbf{1}, \epsilon$ of $\mu_2$.
            The representation $\pi_{\gen}[\chi]$ is generic, and $\pi_{\mathrm{sc}}[\chi]$ is non-generic.
            \item 
            If $\sigma(W_F) = \mu_3$, then the $L$-packet $\Pi_{\phi}$ consists of three representations $\pi_{\gen}[\chi], \pi_{\mathrm{sc}, 1}[\chi], \pi_{\mathrm{sc}, 2}[\chi]$.
            The representation $\pi_{\gen}[\chi] = \pi_{\gen}[\chi^2]$ is generic, and $\pi_{\mathrm{sc}, 1}[\chi], \pi_{\mathrm{sc}, 2}[\chi]$ are supercuspidal.
            \item 
            If $\sigma(W_F) = S_3$, then the $L$-packet $\Pi_{\phi}$ consists of a single generic representation $\pi_{\gen}[\rho \circ \sigma]$.
      \end{enumerate}

      In Cases $(1), (2), (3)$ all the non-supercuspidal representations in the $L$-packets have the cuspidal support in the maximal torus of ${\rG_2}$.
      In Case $(4)$, the generic representation $\pi_{\gen}[\rho \circ \sigma]$ has the cuspidal support in the proper Levi subgroup $M_Q$ corresponding to the long root of ${\rG_2}$.
\end{theorem}

\begin{proposition}\label{prop:tempered-subregular}
      The stable distribution $S\Theta_{\phi_{\sigma}}$ of the tempered subregular $L$-packet $\Pi_{\phi_{\sigma}}$ is given by
      \begin{align*}
            S\Theta_{\phi_{\sigma}} = \sum_{\pi \in \Pi_{\phi}} m_{\pi} \Theta_{\pi},
      \end{align*}
      where $m_{\pi}$ is the dimension of the corresponding irreducible representation of $\cS_{\phi_{\sigma}} = \Cent_{S_3}(\sigma(W_F))$.

      We have the standard endoscopic character identities:
      \begin{align*}
            \Trans^{\rG_2, *}_{\SO_4}(\Theta_{\St_{\SO_4}}) &= \Theta_{\pi_{\gen}[1]} - \Theta_{\pi_{\mathrm{sc}}[1]}, \\
            \Trans^{\rG_2, *}_{\SO_4}(\Theta_{\St_{\SO_4} \otimes \chi \circ \sn}) &= \Theta_{\pi_{\gen}[\chi]} - \Theta_{\pi_{\mathrm{sc}}[\chi]}, &&(\chi^2=1,\ \chi\ne1) \\
            \Trans^{\rG_2, *}_{\PGL_3}(\Theta_{\St_{\PGL_3}}) &= \Theta_{\pi_{\gen}[1]} - \Theta_{\pi_{\deg}[1]} + \Theta_{\pi_{\mathrm{sc}}[1]}.
      \end{align*}

      For a non-trivial cubic character $\chi$, we can choose a primitive root of unity $\zeta$ so that we have 
      \begin{align*}
            \Trans^{\rG_2, *}_{\PGL_3}(\Theta_{\St_{\PGL_3} \otimes \chi \circ \det}) &= 
            \Theta_{\pi_{\gen}[\chi]} +\zeta \Theta_{\pi_{\mathrm{sc}, 1}[\chi]} + \zeta^2 \Theta_{\pi_{\mathrm{sc}, 2}[\chi]}, \\
            \Trans^{\rG_2, *}_{\PGL_3}(\Theta_{\St_{\PGL_3} \otimes \chi^2 \circ \det}) &= \Theta_{\pi_{\gen}[\chi]} + \zeta^2 \Theta_{\pi_{\mathrm{sc}, 1}[\chi]} + \zeta \Theta_{\pi_{\mathrm{sc}, 2}[\chi]}.
      \end{align*} 
\end{proposition}

\subsection{Aubert dual of tempered packets}\label{sec:subregular-aubert}

\begin{proposition}\label{prop:subregular-abelian-Aubert-dual-computation}
      In $R(\rG_2)$, we have the following results on the Aubert duals of the tempered subregular $L$-packets of ${\rG_2}$.

      \begin{enumerate}
      \item 
      We have 
      \begin{align*}
            \widehat{\pi_{\gen}[1]} &= J_Q(\pi(1, 1) \nu_2), \\
            \widehat{\pi_{\deg}[1]} &= J_P(\St_{\GL_2}\nu_2^{1/2}), \\
            \widehat{J_Q(\St_{\GL_2}\nu_2^{1/2})} &= J_Q(\St_{\GL_2}\nu_2^{1/2}).
      \end{align*}
      Furthermore, 
      \begin{align*}
            [I_P(\mathbf{1}_{\GL_2}\nu_2^{1/2})] = [\pi_{\deg}[1]] + [\widehat{\pi_{\gen}[1]}] + [J_Q(\St_{\GL_2}\nu_2^{1/2})], \\
            [I_P(\St_{\GL_2}\nu_2^{1/2})] 
            =
             [\pi_{\gen}[1]] + [\widehat{\pi_{\deg}[1]}] + [J_Q(\St_{\GL_2}\nu_2^{1/2})], \\
            [I_Q(\St_{\GL_2}\nu_2^{1/2})] = [\pi_{\gen}[1]] + [\pi_{\deg}[1]] + [J_Q(\St_{\GL_2}\nu_2^{1/2})],
      \end{align*}
      and
      \begin{align*}
            [J_Q(\pi(1, 1) \nu_2)] &= [I_Q(\pi(1, 1)\nu_2)] - [I_P(\St_{\GL_2}\nu_2^{1/2})] - [I_Q(\St_{\GL_2}\nu_2^{1/2})] + [\pi_{\mathrm{gen}}[1]], \\
            [J_P(\St_{\GL_2}\nu_2^{1/2})] &= [I_P(\St_{\GL_2}\nu_2^{1/2})] - [I_Q(\St_{\GL_2}\nu_2^{1/2})] + [\pi_{\deg}[1]]. \\
      \end{align*}

      \item
      Let $\chi$ be a non-trivial quadratic character of $F^{\times}$.
      We have
      \begin{align*}
            \widehat{\pi_{\gen}[\chi]} &= J_Q(\pi(1, \chi) \nu_2), \\
            \widehat{J_P(\St_2(\chi)\nu_2^{1/2})} &= J_Q(\St_2(\chi)\nu_2^{1/2}).
      \end{align*}
      Furthermore, 
      \begin{align*}
            [I_P(\chi \circ \det \nu_2^{1/2})] = [\widehat{\pi_{\gen}[\chi]}] + [J_Q(\St_2(\chi)\nu_2^{1/2})], \\
            [I_P(\St_2(\chi)\nu_2^{1/2})] = [\pi_{\gen}[\chi]] + [J_P(\St_2(\chi)\nu_2^{1/2})], \\
            [I_Q(\St_2(\chi)\nu_2^{1/2})] = [\pi_{\gen}[\chi]] + [J_Q(\St_2(\chi)\nu_2^{1/2})].
      \end{align*}
      and
      \begin{align*}
            [J_Q(\pi(1, \chi) \nu_2)] &= [I_Q(\pi(1, \chi)\nu_2)] - [I_P(\St_2(\chi)\nu_2^{1/2})] - [I_Q(\St_{2}(\chi)\nu_2^{1/2})] + [\pi_{\mathrm{gen}}[\chi]].
      \end{align*}

      \item
      Let $\chi$ be a non-trivial cubic character of $F^{\times}$.
      We have
      \begin{align*}
            \widehat{\pi_{\gen}[\chi]} &= J_Q(\pi(\chi, \chi^{-1}) \nu_2), \\
            \widehat{J_P(\St_2(\chi)\nu_2^{1/2})} &= J_P(\St_2(\chi^{-1})\nu_2^{1/2}).
      \end{align*}
      Furthermore, 
      \begin{align*}
            [I_P(\chi \circ \det \nu_2^{1/2})] &= 
            [\widehat{\pi_{\gen}[\chi]}] + [J_P(\St_2(\chi^{-1})\nu_2^{1/2})], \\
            [I_P(\St_2(\chi^{-1})\nu_2^{1/2})] &=
            [\pi_{\gen}[\chi]] + [J_P(\St_2(\chi^{-1})\nu_2^{1/2})],
      \end{align*}
      and we have
      \begin{align*}
            [J_Q(\pi(\chi, \chi^{-1}) \nu_2)]
             &= [I_Q(\pi(\chi, \chi^{-1})\nu_2)] - [I_P(\St_2(\chi)\nu_2^{1/2})] \\
             &- 
             [I_P(\St_{2}(\chi^{-1})\nu_2^{1/2})] + [\pi_{\mathrm{gen}}[\chi]].
      \end{align*}
      \end{enumerate}
\end{proposition}

\begin{proof}
      The composition series are due to Mui\'c
      \cite{Muic1997-UnitaryDualG2}*{Propositions 4.3, 4.1, and 4.2}.
      We use their formulations in
      \cite{Qin2026-AubertZelevinskyInvolutionG2HeckeAlgebras}*{
      Propositions 6.2.1(2), 7.4.1(2), and 7.5.1(2)}, with the
      inducing data expressed in the intrinsic $\GL_2$ coordinates of
      Section~\ref{subsec:Qin-notation}.
      In the first case Qin's $\pi'(1)$ is $\pi_{\gen}[1]$ and $\pi(1)$ is $\pi_{\deg}[1]$;
      the identification also follows from the generic and nongeneric
      subrepresentations in
      \cite{GanSavin2023-HoweDualityDichotomyExceptionalThetaCorrespondences}*{
      Propositions 3.1 and 3.2}.
      The corresponding duality assertions are
      \cite{Qin2026-AubertZelevinskyInvolutionG2HeckeAlgebras}*{
      Propositions 6.2.3, 7.4.3, and 7.5.3}. Qin uses
      $\sum_{I\subset\Delta}(-1)^{|I|}I_{P_I}r_{P_I}$, which equals
      \eqref{eq:Aubert-operator} since $\operatorname{rank}{\rG_2}=2$.
      All these principal-series constituents have toral support, so their
      support sign is $+1$. Thus these duality assertions agree with positive
      Aubert duality. The supercuspidal case follows from
      Proposition~\ref{prop:Aubert}.
\end{proof}


\begin{proposition}
      Assume that $\sigma(W_F) = S_3$.
      Let $\tau$ be a supercuspidal representation of $\GL_2(F)$ corresponding to the $L$-parameter $\rho \circ \sigma \colon W_F \to \GL_2(\C)$.
      We have an exact sequence
      \begin{align*}
            0 \to \pi_{\gen}[\tau] \to I_Q(\tau \nu_2) \to J_Q(\tau \nu_2) \to 0.
      \end{align*}

      We also have 
      \begin{align*}
            \widehat{\pi_{\gen}[\tau]} &= J_Q(\tau \nu_2).
      \end{align*}
\end{proposition}

\begin{proof}
      This follows from \cite{GanSavin2023-HoweDualityDichotomyExceptionalThetaCorrespondences}*{Proposition 3.2 (i), the second bullet point}.
\end{proof}


\subsection{The case with trivial image}\label{sec:subregular-trivial}

We consider the case where $\sigma(W_F) = \{1\}$.
In this case, we have $\cS_{\psi} = S_3$.
Let $\psi'$ (resp. $\phi'$) be the lift of $\psi$ (resp. $\phi$) to $\PGSO_8(F)$.

\begin{proposition}\label{prop:subregular-trivial-A-packet-lift-PGSO8}
      The packet $\Pi_{\psi'}$ consists of two representations of $\PGSO_8(F)$
      \begin{align*}
            J_{Q'}((\pi(1, 1)\nu_2)_{M_{Q'}})
            \quad
            \text{and}
            \quad 
            J_{P'} ((\St_{\GL_2}\nu_2^{1/2})_{M_{P'}}).
      \end{align*}
\end{proposition}

\begin{proof}
      This follows from \cite{BanJantzen2003-DegeneratePrincipalSeriesEvenOrthogonalGroups}*{Theorem 5.3, (3)(b)} with $k=3, l=1, n=0, \rho = \mathbf{1}, \sigma = \mathbf{1}$ and $\alpha = 0$ and Proposition \ref{prop:Xu-Aubert} since the cuspidal support of these representations coincides with the cuspidal support of the lift of $\phi$.
\end{proof}

\begin{proposition}\label{prop:subregular-trivial-ecr}
      The distribution
      \begin{align*}
            S\Theta_{\psi} := \Theta_{J_Q(\pi(1, 1) \nu_2)} + 2\Theta_{J_P(\St_{\GL_2}\nu_2^{1/2})} + \Theta_{\pi_{\mathrm{sc}}[1]}
      \end{align*}
      is equal to 
      \begin{align*}
            \Theta_{I_Q(\pi(1, 1) \nu_2)} + \Theta_{I_P(\St_{\GL_2}\nu_2^{1/2})} -3 \Theta_{I_Q(\St_{\GL_2}\nu_2^{1/2})} + S\Theta_{\phi}
      \end{align*}
      and is stable. This transfers to the distribution
      \begin{align*}
            \cT S\Theta_{\psi} =
            \Theta_{\widetilde{J_{Q'}((\pi(1, 1)\nu_2)_{M_{Q'}})}} 
            +
            \Theta_{\widetilde{J_{P'} ((\St_{\GL_2}\nu_2^{1/2})_{M_{P'}})}}.
      \end{align*}

      We set 
      \begin{align*}
            \pi_{\psi}(\mathbf{1}) &= J_Q(\pi(1, 1) \nu_2) \\
            \pi_{\psi}(\rho) &= J_P(\St_{\GL_2}\nu_2^{1/2}) \\
            \pi_{\psi}(\epsilon) &= \pi_{\mathrm{sc}}[1].
      \end{align*}
      Then, we have 
      \begin{align*}
            \Trans^{\rG_2, *}_{\SO_4}(\Theta_{\mathbf{1}_{\SO_4}}) &= \Theta_{\pi_{\psi}(\mathbf{1})} - \Theta_{\pi_{\psi}(\epsilon)}, \\
            \Trans^{\rG_2, *}_{\PGL_3}(\Theta_{\mathbf{1}_{\PGL_3}}) &= \Theta_{\pi_{\psi}(\mathbf{1})} - \Theta_{\pi_{\psi}(\rho)} + \Theta_{\pi_{\psi}(\epsilon)}.
      \end{align*}
\end{proposition}

\begin{proof}
      This follows from Proposition \ref{prop:subregular-abelian-Aubert-dual-computation} (1), Proposition \ref{prop:subregular-trivial-A-packet-lift-PGSO8} and Theorem \ref{thm:signed-transfer}.
\end{proof}

\subsection{The cases with quadratic and cubic image}\label{sec:subregular-quadratic}

We now consider the case where $\sigma(W_F) = \mu_2$ or $\mu_3$.

\begin{proposition}\label{prop:subregular-quadratic-A-packet-lift-SO8}
      If $\chi$ is a non-trivial quadratic character of $F^{\times}$, then the packet $\Pi_{\phi'}$ of $\PGSO_8(F)$ consists of discrete series representations $\pi'_{\gen}[\chi]$ and $\pi'_{\deg}[\chi]$ corresponding to the two irreducible representations $\mathbf{1}$ and $\epsilon$ of $\mu_2$.
      The latter is supercuspidal.
      Thus, the packet $\Pi_{\psi'}$ of $\PGSO_8(F)$ consists of 
      \begin{align*}
            J_{Q'}
            ((\pi(1, \chi)\nu_2)_{M_{Q'}}) \quad \text{and} \quad \pi'_{\deg}[\chi].
      \end{align*}
\end{proposition}

\begin{proof}
      This follows from Lemma \ref{lem:triality-invariant-non-tempered}, \cite{Xu2017-MoeglinParametrizationArthurPackets}*{Definition 6.3(i), (iii)(c)(1), p.920} and \cite{Xu2017-CuspidalSupportDiscreteSeries}*{Theorem 3.3}.
\end{proof}


\begin{proposition}\label{prop:subregular-cubic-A-packet-lift-SO8}
      Let $\chi$ be a non-trivial cubic character of $F^{\times}$.
      Then, the packet $\Pi_{\phi'}$ of $\PGSO_8(F)$ is a singleton consisting of the representation
      \begin{align*}
            I^{\PGSO_8}_{\GL_3 \times \GL_1 \times \GL_1/\GL_1}(\St_{3}(\chi) \boxtimes \mathbf{1} \boxtimes \mathbf{1}). 
      \end{align*} 
      Thus, the packet $\Pi_{\psi'}$ of $\PGSO_8(F)$ is also a singleton consisting of the representation
      \begin{align*}
            I^{\PGSO_8}_{\GL_3 \times \GSO_2 / \GL_1}(\chi \circ \det \boxtimes \mathbf{1}_{\GSO_2}) 
            = I^{\PGSO_8}_{\GL_3 \times \GSO_2 / \GL_1}(\chi^{-1} \circ \det \boxtimes \mathbf{1}_{\GSO_2}).
      \end{align*}
      This representation is also equal to 
      \begin{align*}
            J^{\PGSO_8}_{Q'}(((\pi(\chi, \chi^{-1})\nu_2)_{M_{Q'}})).
      \end{align*}
\end{proposition}

\begin{proof}
      The first assertion follows from the fact that the standard parameter of $\phi'$ is given as
      \begin{align*}
            \chi \boxtimes S_3^D \oplus \chi^{-1} \boxtimes S_3^D \oplus 2 \cdot \mathbf{1} 
      \end{align*}
      The second assertion follows from Proposition \ref{prop:Aubert}.
      The last assertion follows from \cite{Xu2024-ArthurPacketsQuasisplitGSp2nGO2n}*{Lemma 9.3} and Lemma \ref{lem:triality-invariant-non-tempered}.
\end{proof}

Then, we obtain the following results from Theorem \ref{thm:signed-transfer} as in the previous subsection for the case with trivial image.

\begin{proposition}\label{prop:subregular-quadratic-ecr}
      For a non-trivial quadratic character $\chi$, the distribution
      \begin{align*}
           S\Theta_{\psi} = \Theta_{J_Q(\pi(1, \chi) \nu_2)} + \Theta_{\pi_{\mathrm{sc}}[\chi]}
      \end{align*}
      is stable and equal to 
      \begin{align*}
            \Theta_{I_Q(\pi(1, \chi) \nu_2)} - \Theta_{I_P(\St_2(\chi)\nu_2^{1/2})} - \Theta_{I_Q(\St_2(\chi)\nu_2^{1/2})} + S\Theta_{\phi}.
      \end{align*}
      This distribution transfers to the distribution
      \begin{align*}
            \cT S\Theta_{\psi} 
            = 
            \Theta_{\widetilde{
                  J_{Q'}
            ((\pi(1, \chi)\nu_2)_{M_{Q'}})
            }} 
            +     
            \Theta_{\widetilde{\pi'_{\deg}[\chi]}}.
      \end{align*}
      Furthermore, we have 
      \begin{align*}
            \Theta_{\widetilde{
             J^{\PGSO_8}_{Q'}(((\pi(1, \chi)\nu_2)_{M_{Q'}}))
            }} 
            -
            \Theta_{\widetilde{\pi'_{\deg}[\chi]}}
            = 
            \Trans^{\widetilde{\PGSO_8}, *}_{\SO_4}(\Theta_{\chi \circ \sn}).
      \end{align*}

      We set 
      \begin{align*}
            \pi_{\psi}(\mathbf{1}) &= J_Q(\pi(1, \chi) \nu_2) \\
            \pi_{\psi}(\epsilon) &= \pi_{\mathrm{sc}}[\chi].
      \end{align*}
      Then, we have
      \begin{align*}
            \Trans^{\rG_2, *}_{\SO_4}(\Theta_{\chi \circ \sn}) &=
            \Theta_{\pi_{\psi}(\mathbf{1})} - \Theta_{\pi_{\psi}(\epsilon)}.
      \end{align*}
\end{proposition}

\begin{proposition}\label{prop:subregular-cubic-ecr}
      For a non-trivial cubic character $\chi$, the distribution
      \begin{align*}
            S\Theta_{\psi} = \Theta_{J_Q(\pi(\chi, \chi^{-1}) \nu_2)} + \Theta_{\pi_{\mathrm{sc}, 1}[\chi]} + \Theta_{\pi_{\mathrm{sc}, 2}[\chi]}
      \end{align*}
      is stable and equal to
      \begin{align*}
            \Theta_{I_Q(\pi(\chi, \chi^{-1}) \nu_2)} - \Theta_{I_P(\St_2(\chi)\nu_2^{1/2})} - \Theta_{I_P(\St_2(\chi^{-1})\nu_2^{1/2})} + S\Theta_{\phi}.
      \end{align*}
      This distribution transfers to the distribution
      \begin{align*}
            \cT S\Theta_{\psi}
            =
            \Theta_{\widetilde{
                 J^{\PGSO_8}_{Q'}(((\pi(\chi, \chi^{-1})\nu_2)_{M_{Q'}}))
            }}.
      \end{align*}

      We set 
      \begin{align*}
            \pi_{\psi}(\mathbf{1}) &= J_Q(\pi(\chi, \chi^{-1}) \nu_2) \\
            \pi_{\psi}(\chi) &= \pi_{\mathrm{sc}, 1}[\chi] \\
            \pi_{\psi}(\chi^2) &= \pi_{\mathrm{sc}, 2}[\chi].
      \end{align*}
      Then, we have
      \begin{align*}
            \Trans^{\rG_2, *}_{\PGL_3}(\Theta_{\chi \circ \det}) &= \Theta_{\pi_{\psi}(\mathbf{1})} + \zeta \Theta_{\pi_{\psi}(\chi)} + \zeta^2 \Theta_{\pi_{\psi}(\chi^2)}, \\
            \Trans^{\rG_2, *}_{\PGL_3}(\Theta_{\chi^2 \circ \det}) &= \Theta_{\pi_{\psi}(\mathbf{1})} + \zeta^2 \Theta_{\pi_{\psi}(\chi)} + \zeta \Theta_{\pi_{\psi}(\chi^2)}.
      \end{align*}
\end{proposition}

\subsection{The case with $S_3$-image}\label{sec:subregular-full}

Let $\tau$ be the supercuspidal representation of $\GL_2(F)$ corresponding to the $L$-parameter $\rho \circ \sigma \colon W_F \to \GL_2(\C)$.
Let $\omega$ be the central character of $\tau$.


Then, we have the following result.
\begin{proposition}\label{prop:subregular-full-ecr}
      The packet $\Pi_{\psi'}$ of $\PGSO_8(F)$ is a singleton consisting of the representation
      \begin{align*}
            J_{Q'}((\tau \nu_2)_{M_{Q'}}).
      \end{align*}

      The distribution 
      \begin{align*}
            \Theta_{J_{Q}(\tau \nu_2)}
      \end{align*}
      is stable and transfers to the distribution
      \begin{align*}
            \cT \Theta_{J_{Q}(\tau \nu_2)} = \Theta_{\widetilde{J_{Q'}((\tau \nu_2)_{M_{Q'}})}}.  
      \end{align*}

      We set 
      \begin{align*}
            \pi_{\psi}(\mathbf{1}) = J_{Q}(\tau \nu_2).
      \end{align*}
\end{proposition}

\begin{proof}
      This follows from Theorem \ref{thm:signed-transfer} with the support sign $-1$ on both sides.
\end{proof}


\section{Co-tempered long-root and short-root parameters: supercuspidal cases}\label{sec:co-tempered-long-short-supercuspidal}

We now consider the co-tempered long-root and short-root parameters of ${\rG_2}$.
This is given by a supercuspidal irreducible representation $\sigma$ of $\PGL_2(F)$.
The swapped parameter $\phi = \widehat\psi$ is tempered with $\cS_{\phi} \cong \mu_2$.
Thus, the $L$-packet $\Pi_{\phi}$ consists of a generic representation $\delta = \delta_{R}(\sigma)$ in Lemma \ref{lem:rank-one} for $R \in \{P, Q\}$ and $s = 1/2$ (see also \cite{GanSavin2023-HoweDualityDichotomyExceptionalThetaCorrespondences}*{Propositions 3.1 and 3.2}) and a non-generic supercuspidal representation $\pi_{\mathrm{sc}}[\sigma]$ by \cite{Xu2017-CuspidalSupportDiscreteSeries}*{Theorem 3.3} and \cite{GanSavin2023-LocalLanglandsConjectureG2}*{Proposition 8.2}.

Furthermore, the $L$-packet attached to the lift of $\phi$ to $\PGSO_8(F)$ consists of two discrete series representations, namely a generic representation $\delta'$ and a non-generic supercuspidal representation $\pi'_{\mathrm{sc}}[\sigma]$.
These representations are obtained from the composition of the exceptional theta lift from $\rG_2$ to $\PGSp_6$ and the similitude theta lift from $\PGSp_6$ to $\PGSO_8$.

We can compute the Aubert dual of $\delta$ in Lemma \ref{lem:rank-one} and obtain the endoscopic character identities by Theorem \ref{thm:signed-transfer}.


\begin{proposition}\label{prop:co-tempered-long-short-ecr}
      We set 
      \begin{align*}
            \pi_{\psi}(\mathbf{1}) &= \widehat{\delta}, \\
            \pi_{\psi}(\epsilon) &= \pi_{\mathrm{sc}}[\sigma].
      \end{align*}
      Then, the distribution
      \begin{align*}
            S\Theta_{\psi} = \Theta_{\pi_{\psi}(\mathbf{1})} - \Theta_{\pi_{\psi}(\epsilon)}
      \end{align*}
      is stable and transfers to the distribution
      \begin{align*}
            \cT S\Theta_{\psi} 
            =
            \Theta_{\widetilde{\widehat{\delta'}}} - \Theta_{\widetilde{\pi'_{\mathrm{sc}}[\sigma]}}.
      \end{align*}

      If $\psi$ is a long root (resp. short root) parameter, then we have
      \begin{align*}
            \delta = \delta_{Q}(\sigma) \quad (\text{resp. } \delta = \delta_{P}(\sigma)),
      \end{align*}
      in \cite{GanSavin2023-HoweDualityDichotomyExceptionalThetaCorrespondences}.

      Furthermore, we have 
      \begin{align*}
            \Theta_{\widetilde{\widehat{\delta'}}} + \Theta_{\widetilde{\pi'_{\mathrm{sc}}[\sigma]}}
            =
            \Trans^{\widetilde{\PGSO_8}, *}_{\SO_4}(\Theta_{J^{\SO_4}_{R}(\sigma \nu_2^{1/2})}), \\
            \Theta_{\pi_{\psi}(\mathbf{1})} + \Theta_{\pi_{\psi}(\epsilon)}
            =
            \Trans^{\rG_2, *}_{\SO_4}(\Theta_{J^{\SO_4}_{R}(\sigma \nu_2^{1/2})}),
      \end{align*}
      where $R = P_{l}$ (resp. $R = P_{s}$) if $\psi$ is a long root (resp. short root) parameter and $J^{\SO_4}_{R}(\sigma \nu_2^{1/2})$ is the Langlands quotient.
\end{proposition}


\section{A-parameters with unbounded image}\label{sec:A-parameter-unbounded-ecr}

In the long-root and short-root cases, we can consider the A-parameters with unbounded image associated with generic unitary principal series representations of $\PGL_2(F)$.
Let $\sigma$ be a generic unitary principal series irreducible representation of $\PGL_2(F)$.
We consider the long-root and short-root A-parameter $\psi$ of $\rG_2(F)$ associated with $\sigma$.
The long-root parameter is of the form
\begin{align*}
      \psi \colon W_F \times \SL_2^A(\C) \xrightarrow{\phi_{\sigma} \times \id_{\SL_2}} \SL_{2, s}(\C) \times \SL_{2, l}(\C) \to \rG_2(\C).
\end{align*}
For the short-root parameter, interchange the subscripts $s$ and $l$
in this formula.
Here, we relax the boundedness condition on the image of $\phi_{\sigma}$.

We first consider the long-root case.
Then, the parameter factors through the long root Levi subgroup of $\rG_2(\C)$ and defines an A-parameter for $\GL_2(F)$.
This corresponds to a character $\chi_{M_P}$ of $M_P$.
In this case, we set 
\begin{align*}
      \pi_{\psi}(\mathbf{1}) = I_P(\chi_{M_P}).
\end{align*}
By \cite{Muic1997-UnitaryDualG2}*{Theorem 3.1 (i), Lemma 5.2} the representation $I_P(\chi_{M_P})$ is irreducible and unitary.
We have the twisted endoscopic character identity
\begin{align*}
      \Trans^{\widetilde{\PGSO_8}, *}_{\rG_2} \Theta_{I_P(\chi_{M_P})} = \Theta_{\widetilde{I_{P'}(\chi_{M_{P'}})}},
\end{align*}
where $\chi_{M_{P'}}$ is the character of $M_{P'}(F)$ corresponding to $\chi_{M_P}$.

The short-root case is similar:
The parameter factors through the short root Levi subgroup of $\rG_2(\C)$ and defines an A-parameter for $\GL_2(F)$.
This corresponds to a character $\chi_{M_Q}$ of $M_Q$.
In this case, we set
\begin{align*}
      \pi_{\psi}(\mathbf{1}) = I_Q(\chi_{M_Q}).
\end{align*}
By \cite{Muic1997-UnitaryDualG2}*{Theorem 3.1 (ii), Lemma 5.2} the representation $I_Q(\chi_{M_Q})$ is irreducible and unitary.
We have the twisted endoscopic character identity
\begin{align*}
      \Trans^{\widetilde{\PGSO_8}, *}_{\rG_2} \Theta_{I_Q(\chi_{M_Q})} = \Theta_{\widetilde{I_{Q'}(\chi_{M_{Q'}})}},
\end{align*}
where $\chi_{M_{Q'}}$ is the character of $M_{Q'}(F)$ corresponding to $\chi_{M_Q}$.

We summarize the above discussion in the following proposition.
\begin{proposition}\label{prop:co-tempered-long-short-unbounded-ecr}
      Let $\psi$ be a long-root (resp. short-root) A-parameter of $\rG_2(F)$ associated with a generic unitary principal series representation $\sigma$ of $\PGL_2(F)$.
      Then, the A-packet $\Pi_{\psi}$ consists of a single representation $\pi_{\psi}(\mathbf{1})$ given by $I_P(\chi_{M_P})$ (resp. $I_Q(\chi_{M_Q})$).
      The distribution $\Theta_{\pi_{\psi}(\mathbf{1})}$ is stable and transfers to the distribution $\Theta_{\widetilde{I_{P'}(\chi_{M_{P'}})}}$ (resp. $\Theta_{\widetilde{I_{Q'}(\chi_{M_{Q'}})}}$).
\end{proposition}


\section{Mixed long root parameters}\label{sec:mixed-long}

Let $\psi_D$ be a mixed long root parameter of ${\rG_2}$ associated with a quadratic character $\chi$ of $F^{\times}$.
Let $\psi'_D$ denote the lift to $\PGSO_8(F)$.


\begin{lemma}  
      There is a triality-invariant A-packet corresponding to $\psi'_D$ for $\PGSO_8(F)$ consisting of two representations.
      One is non-tempered and the other is tempered.

      If $\chi$ is trivial, then the non-tempered representation is given by
      \begin{align*}
            J_{Q'}((\St_{\GL_2} \nu_2^{1/2})_{M_{Q'}}),
      \end{align*}
      and the tempered representation has the standard parameter
      \begin{align*}
            2 \cdot S_3^D \oplus 2 \cdot \mathbf{1}
      \end{align*}
      and it is non-generic.

      If $\chi$ is non-trivial, then the non-tempered representation is given by
      \begin{align*}
            J_{Q'}((\St_{2}(\chi) \nu_2^{1/2})_{M_{Q'}}),
      \end{align*}
      and the other representation is supercuspidal and non-generic with the standard parameter 
      \begin{align*}
            \chi \boxtimes S_3^D \oplus \chi \oplus S_3^D \oplus \mathbf{1}.
      \end{align*}
\end{lemma}

\begin{proof}
      We first determine the A-packet for $\SO_8(F)$ corresponding to the lift of $\psi_D$.
      If $\chi$ is trivial, then \cite{Xu2021-NonarchimedeanComponentsNonEndoscopicAutomorphicRepresentations}*{Theorem 1.3 (2)} implies the A-packet consists of two representations, one is non-tempered and given by
      \begin{align*}
            J_{\GL_2 \times \SO_4}(\St_{\GL_2} \nu_2^{1/2} \boxtimes \St_{\SO_4})
      \end{align*}
      and the other is a tempered and non-generic representation $\pi'_{\temp, \mathrm{ng}}$ with the standard parameter $2 \cdot S_3^D \oplus 2 \cdot \mathbf{1}$.
      If $\chi$ is non-trivial, then we can use \cite{Xu2017-MoeglinParametrizationArthurPackets}*{Theorem 7.8}, which implies that the A-packet for $\SO_8(F)$ consists of three representations, one is non-tempered and given by
      \begin{align*}
            J_{\GL_2 \times \SO_4}(\St_{2}(\chi) \nu_2^{1/2} \boxtimes \St_{\SO_4})
      \end{align*}
      and the others are supercuspidal and non-generic with the standard parameter $\chi \boxtimes S_3^D \oplus \chi \oplus S_3^D \oplus \mathbf{1}$.
      By \cite{Xu2018-LPacketsQuasisplitGSp2nGO2n}*{Lemma 6.6}, these supercuspidal representations form a single $\PGSO_8(F)$-orbit $\pi'_{\mathrm{sc}, \chi}$.
      
      These observations and Lemma \ref{lem:triality-invariant-non-tempered} imply that if there is an A-packet with the standard parameter $\Std_{\PGSO_8} \circ \psi'_D$ and triality-invariant, then it is the A-packet for $\PGSO_8(F)$ containing 
      \begin{align*}
             J_{Q'}((\St_{\GL_2} \nu_2^{1/2})_{M_{Q'}})
      \end{align*}
      if $\chi$ is trivial and
      \begin{align*}
             J_{Q'}((\St_{2}(\chi) \nu_2^{1/2})_{M_{Q'}})
      \end{align*}
      if $\chi$ is non-trivial.
      We now show that they actually form an A-packet for $\PGSO_8(F)$.
      By \cite{Xu2024-ArthurPacketsQuasisplitGSp2nGO2n}*{Theorem 8.1 (1)}, it suffices to show the stability of the distribution
      \begin{align*}
            \Theta_{J_{Q'}((\St_{\GL_2} \nu_2^{1/2})_{M_{Q'}})} - \Theta_{\pi'_{\temp, \mathrm{ng}}}, \\
            \Theta_{J_{Q'}((\St_{2}(\chi) \nu_2^{1/2})_{M_{Q'}})} - \Theta_{\pi'_{\mathrm{sc}, \chi}}.
      \end{align*}
      For $\chi = \mathbf{1}$, we can apply \cite{Muic2005-ReducibilityGeneralizedPrincipalSeries}*{Lemma 5.2} to see that the parabolic induction $I_{Q'}((\St_{\GL_2} \nu_2^{1/2})_{M_{Q'}})$ is of length two and has a unique irreducible tempered generic subrepresentation $\pi'_{\temp, \gen}$.
      The standard parameter of this representation must be $2 \cdot S_3^D \oplus 2 \cdot \mathbf{1}$.
      Thus, we have 
      \begin{align*}
            [J_{Q'}((\St_{\GL_2} \nu_2^{1/2})_{M_{Q'}})] - [\pi'_{\temp, \mathrm{ng}}]
            = 
            [I_{Q'}((\St_{\GL_2} \nu_2^{1/2})_{M_{Q'}})] - [\pi'_{\temp, \gen}] - [\pi'_{\temp, \mathrm{ng}}].
      \end{align*}
      The corresponding distribution is stable by Theorem \ref{thm:tempered-input}.

      The case for non-trivial $\chi$ is similar by using \cite{Muic2005-ReducibilityStandardRepresentations}*{Lemma 6.1}.
      In this case, we obtain the exact sequence
      \begin{align*}
            0 \to \pi'_{\gen, \chi} \to I_{Q'}((\St_{2}(\chi) \nu_2^{1/2})_{M_{Q'}}) \to J_{Q'}((\St_{2}(\chi) \nu_2^{1/2})_{M_{Q'}}) \to 0,
      \end{align*}
      where $\pi'_{\gen, \chi}$ is the unique irreducible generic discrete series subrepresentation of $I_{Q'}((\St_{2}(\chi) \nu_2^{1/2})_{M_{Q'}})$ with the standard parameter $\chi \boxtimes S_3^D \oplus \chi \oplus S_3^D \oplus \mathbf{1}$.
      Thus, we have 
      \begin{align*}
            [J_{Q'}((\St_{2}(\chi) \nu_2^{1/2})_{M_{Q'}})] - [\pi'_{\mathrm{sc}, \chi}] 
            =
            [I_{Q'}((\St_{2}(\chi) \nu_2^{1/2})_{M_{Q'}})] - ([\pi'_{\gen, \chi}] + [\pi'_{\mathrm{sc}, \chi}]).
      \end{align*}
      The corresponding distribution is stable by Theorem \ref{thm:tempered-input}.
\end{proof}

We record the following corollary of the proof above.
\begin{corollary}
      We have the following exact sequences
      \begin{align*}
            0 \to \pi'_{\temp, \gen} \to I_{Q'}((\St_{\GL_2} \nu_2^{1/2})_{M_{Q'}}) \to J_{Q'}((\St_{\GL_2} \nu_2^{1/2})_{M_{Q'}}) \to 0, \\
            0 \to \pi'_{\mathrm{gen}, \chi} \to I_{Q'}((\St_{2}(\chi) \nu_2^{1/2})_{M_{Q'}}) \to J_{Q'}((\St_{2}(\chi) \nu_2^{1/2})_{M_{Q'}}) \to 0.
      \end{align*}
\end{corollary}

\begin{proposition}\label{prop:mixed-long-root-ecr}
      Let $\psi_D$ be a mixed long root parameter of ${\rG_2}$ associated with a quadratic character $\chi$ of $F^{\times}$.
      \begin{enumerate}
            \item 
            If $\chi$ is trivial, then the distribution
            \begin{align*}
                  S\Theta_{\psi_D} = \Theta_{J_Q(\St_{\GL_2} \nu_2^{1/2})} - (\Theta_{\pi_{\deg}[1]} + \Theta_{\pi_{\mathrm{sc}}[1]}),
            \end{align*}
            is stable and it is equal to
            \begin{align*}
                  \Theta_{I_Q(\St_{\GL_2} \nu_2^{1/2})} - (\Theta_{\pi_{\gen}[1]} + 2\Theta_{\pi_{\deg}[1]} + \Theta_{\pi_{\mathrm{sc}}[1]}).
            \end{align*}
            This distribution transfers to the distribution
            \begin{align*}
                  \cT S\Theta_{\psi_D} =
                  \Theta_{\widetilde{J_{Q'}((\St_{\GL_2} \nu_2^{1/2})_{M_{Q'}})}} - \Theta_{\widetilde{\pi'_{\temp, \mathrm{ng}}}},
            \end{align*}
            where $\pi'_{\temp, \mathrm{ng}}$ is the triality-invariant tempered non-generic representation of $\PGSO_8(F)$ with the standard parameter $2 \cdot S_3^D \oplus 2 \cdot \mathbf{1}$.

            Furthermore, we have distributions
            \begin{align*}
                  \Theta_{J_Q(\St_{\GL_2} \nu_2^{1/2})} + (\Theta_{\pi_{\deg}[1]} + \Theta_{\pi_{\mathrm{sc}}[1]})
            \end{align*}
            and 
            \begin{align*}
                  \Theta_{\widetilde{J_{Q'}((\St_{\GL_2} \nu_2^{1/2})_{M_{Q'}})}} + \Theta_{\widetilde{\pi'_{\temp, \mathrm{ng}}}},
            \end{align*}
            which are transfers of the stable distribution
            \begin{align*}
                  \Theta_{J_{P_{l}}^{\SO_4}(\St_{\GL_2} \nu_2^{1/2})}
            \end{align*}
            of $\SO_4(F)$.
            Here, we recall that $P_{l}$ is the long root parabolic subgroup of $\SO_4(F)$ with the short root dual parabolic subgroup.

            We set 
            \begin{align*}
                  \pi_{\psi_D}(\mathbf{1}) &= J_Q(\St_{\GL_2} \nu_2^{1/2}), \\
                  \pi_{\psi_D}(\epsilon) &= \pi_{\deg}[1] \oplus \pi_{\mathrm{sc}}[1].
            \end{align*}

            \item 
            If $\chi$ is non-trivial, then the distribution
            \begin{align*}
                  S\Theta_{\psi_D} = 
                  \Theta_{J_Q(\St_{2}(\chi) \nu_2^{1/2})} - \Theta_{\pi_{\mathrm{sc}}[\chi]},
            \end{align*}
            is stable and it is equal to
            \begin{align*}
                  \Theta_{I_Q(\St_{2}(\chi) \nu_2^{1/2})} - (\Theta_{\pi_{\gen}[\chi]} + \Theta_{\pi_{\mathrm{sc}}[\chi]}).
            \end{align*}
            This distribution transfers to the distribution
            \begin{align*}
                  \cT S\Theta_{\psi_D} =
                  \Theta_{\widetilde{J_{Q'}((\St_{2}(\chi) \nu_2^{1/2})_{M_{Q'}})}} - \Theta_{\widetilde{\pi'_{\mathrm{sc}, \chi}}},
            \end{align*}
            where $\pi'_{\mathrm{sc}, \chi}$ is the triality-invariant supercuspidal non-generic representation of $\PGSO_8(F)$ with the standard parameter
            \begin{align*}
                  \chi \boxtimes S_3^D \oplus \chi \oplus S_3^D \oplus \mathbf{1}.
            \end{align*}

            Furthermore, we have distributions
            \begin{align*}
                  \Theta_{J_Q(\St_{2}(\chi) \nu_2^{1/2})} + \Theta_{\pi_{\mathrm{sc}}[\chi]}
            \end{align*}
            and
            \begin{align*}
                  \Theta_{\widetilde{J_{Q'}((\St_{2}(\chi) \nu_2^{1/2})_{M_{Q'}})}} + \Theta_{\widetilde{\pi'_{\mathrm{sc}, \chi}}},
            \end{align*}
            which are transfers of the stable distribution
            \begin{align*}
                  \Theta_{J_{P_{l}}^{\SO_4}(\St_{2}(\chi) \nu_2^{1/2})}
            \end{align*}
            of $\SO_4(F)$.

            We set 
            \begin{align*}
                  \pi_{\psi_D}(\mathbf{1}) &= J_Q(\St_{2}(\chi) \nu_2^{1/2}), \\
                  \pi_{\psi_D}(\epsilon) &= \pi_{\mathrm{sc}}[\chi].
            \end{align*}
      \end{enumerate}
\end{proposition}

\begin{proof}
      The stability follows from Theorem \ref{thm:tempered-input} and the transfer follows from Theorem \ref{thm:signed-transfer} and Theorem \ref{thm:tempered-input}.
      In these cases, the component group $\cS_{\psi'_D}$ is of order two and the corresponding endoscopic group is $\SO_4(F)$.
\end{proof}

\section{Mixed short root parameters}\label{sec:mixed-short}

By applying the (twisted) Aubert duality, we can obtain the results for the mixed short root parameters of ${\rG_2}$ from the results for the mixed long root parameters in Proposition \ref{prop:mixed-long-root-ecr}.

\begin{proposition}\label{prop:mixed-short-root-trivial-ecr}
      We have the following results for the mixed short root parameters of ${\rG_2}$.
      \begin{enumerate}
            \item 
            In $R(\rG_2)$, we have 
            $\widehat{J_Q(\St_{\GL_2}\nu_2^{1/2})} = J_Q(\St_{\GL_2}\nu_2^{1/2})$ 
            and
            $\widehat{\pi_{\deg}[1]} = J_P(\St_{\GL_2}\nu_2^{1/2})$.
            \item 
            In $R(\PGSO_8 \rtimes \theta)$, we have 
            \begin{align*}
                  \widetilde{\widehat{J_{Q'}((\St_{\GL_2} \nu_2^{1/2})_{M_{Q'}})}}
                  &= 
                  \widetilde{J_{Q'}((\St_{\GL_2} \nu_2^{1/2})_{M_{Q'}})}, \\
                  \widetilde{\widehat{\pi'_{\temp, \mathrm{ng}}}}
                  &= 
                  \widetilde{J_{P'}((\St_{\GL_2}\nu_2^{1/2})_{M_{P'}})}.
            \end{align*}
      \end{enumerate}
      
      We set 
      \begin{align*}
            \pi_{\psi_A}(\mathbf{1}) &= J_P(\St_{\GL_2}\nu_2^{1/2}) \oplus \pi_{\mathrm{sc}}[1], \\
            \pi_{\psi_A}(\epsilon) &= J_Q(\St_{\GL_2}\nu_2^{1/2}).
      \end{align*}
\end{proposition}

\begin{proof}
      The statement $(1)$ follows from Proposition \ref{prop:subregular-abelian-Aubert-dual-computation}.
      Since the Aubert duality preserves cuspidal supports, the first equation of the statement $(2)$ follows from \cite{BanJantzen2003-DegeneratePrincipalSeriesEvenOrthogonalGroups}*{Theorem 5.3} with $\rho = \mathbf{1}, \sigma_0 = \mathbf{1}, n = 0, k = l = 2$ and $\alpha = -\frac{1}{2}$ and Proposition \ref{prop:Xu-Aubert}.
      Since by \cite{Xu2024-ArthurPacketsQuasisplitGSp2nGO2n}*{Lemma 9.2} and Lemma \ref{lem:triality-invariant-non-tempered}, the representation $J_{P'}((\St_{\GL_2}\nu_2^{1/2})_{M_{P'}})$ must be contained in the A-packet for $\PGSO_8(F)$ corresponding to $\psi'_{A}$, the second equation of the statement $(2)$ follows.
\end{proof}


\begin{proposition}\label{prop:mixed-short-root-quadratic-ecr}
      We have the following results for the mixed short root parameters of ${\rG_2}$ associated with a non-trivial quadratic character $\chi$.
      \begin{enumerate}
            \item 
            In $R(\rG_2)$, we have 
            $\widehat{J_Q(\St_{2}(\chi)\nu_2^{1/2})} = J_P(\St_{2}(\chi)\nu_2^{1/2})$.
            \item 
            In $R(\PGSO_8 \rtimes \theta)$, we have 
            \begin{align*}
                  \widetilde{\widehat{J_{Q'}((\St_{2}(\chi) \nu_2^{1/2})_{M_{Q'}})}}
                  &= 
                  \widetilde{J_{P'}((\St_{2}(\chi) \nu_2^{1/2})_{M_{P'}})}, \\
                  \widetilde{\widehat{\pi'_{\mathrm{sc}, \chi}}} 
                  &=
                  \widetilde{\pi'_{\mathrm{sc}, \chi}}.
            \end{align*}
      \end{enumerate}

      We set
      \begin{align*}
            \pi_{\psi_A}(\mathbf{1}) &= J_P(\St_{2}(\chi)\nu_2^{1/2}), \\
            \pi_{\psi_A}(\epsilon) &= \pi_{\mathrm{sc}}[\chi].
      \end{align*}
\end{proposition}

\begin{proof}
      It suffices to show that the representation $J_{P'}((\St_{2}(\chi) \nu_2^{1/2})_{M_{P'}})$ is contained in the A-packet for $\PGSO_8(F)$ corresponding to $\psi'_{A}$.
      This follows from \cite{Xu2024-ArthurPacketsQuasisplitGSp2nGO2n}*{Lemma 9.3} and Lemma \ref{lem:triality-invariant-non-tempered}.
\end{proof}


\section{Summary}\label{sec:summary}

\begin{definition}
      Let $F$ be a non-archimedean local field of characteristic zero.  
      Let $\psi$ be a non-tempered A-parameter of $\rG_2(F)$.
      \begin{enumerate}
            \item 
            We have defined a map from the set of irreducible representations of the component group $\cS_{\psi}$ to the set of finite length representations of $\rG_2(F)$
            in Propositions \ref{prop:principal-ecr}, \ref{prop:subregular-trivial-ecr}, \ref{prop:subregular-quadratic-ecr}, \ref{prop:subregular-cubic-ecr}, \ref{prop:subregular-full-ecr}, \ref{prop:co-tempered-long-short-ecr}, \ref{prop:mixed-long-root-ecr}, \ref{prop:mixed-short-root-trivial-ecr}, \ref{prop:mixed-short-root-quadratic-ecr} and Proposition \ref{prop:co-tempered-long-short-unbounded-ecr}.
            For $\rho \in \Irr_{\C}(\cS_{\psi})$, let $\pi_{\psi}(\rho)$ denote the corresponding representation of $\rG_2(F)$ of finite length.
            Then, we define the A-packet $\Pi_{\psi}$ of $\rG_2(F)$ to be the set of representations
            \begin{align*}
                  \Pi_{\psi} = \bigcup_{\rho \in \Irr_{\C}(\cS_{\psi})} \mathrm{JH}(\pi_{\psi}(\rho)).
            \end{align*}
            For $s \in \cS_{\psi}$, we also set 
            \begin{align*}
                  \Theta_{\psi, s} = \sum_{\rho \in \Irr_{\C}(\cS_{\psi})} \tr\rho (s s_{\psi}) \Theta_{\pi_{\psi}(\rho)}.
            \end{align*}
            \item 
            Let $\psi'$ be the lift of $\psi$ to $\PGSO_8(F)$.
            Then, we define the A-packet $\Pi_{\psi'}$ of $\PGSO_8(F)$ to be Xu's A-packet appearing in Propositions \ref{prop:principal-ecr}, \ref{prop:subregular-trivial-ecr}, \ref{prop:subregular-quadratic-ecr}, \ref{prop:subregular-cubic-ecr}, \ref{prop:subregular-full-ecr}, \ref{prop:co-tempered-long-short-ecr}, \ref{prop:mixed-long-root-ecr}, \ref{prop:mixed-short-root-trivial-ecr}, \ref{prop:mixed-short-root-quadratic-ecr} and Proposition \ref{prop:co-tempered-long-short-unbounded-ecr}.
            Each A-packet consists of $\Aut(\PGSO_8)(F)$-invariant representations of $\PGSO_8(F)$ and there exists a map
            \begin{align*}
                  \Pi_{\psi'} \to \Irr(\cS_{\psi'}) \colon \pi' \mapsto \langle \cdot, \pi' \rangle_{\psi'},
            \end{align*}
            constructed by Arthur and Xu, see \cite{Xu2024-ArthurPacketsQuasisplitGSp2nGO2n}*{Corollary 6.2}.
      \end{enumerate}
\end{definition}

\begin{remark}
      We use similar notation for $\SO_4$ and $\PGL_3$.
      For an A-parameter $\psi$ of $\SO_4(F)$ or $\PGL_3(F)$, we denote by $\Pi_{\psi}$ the A-packet of $\SO_4(F)$ or $\PGL_3(F)$ and let $S\Theta_{\psi}$ be the stable distribution of $\SO_4(F)$ or $\PGL_3(F)$ associated with $\psi$.
\end{remark}

\begin{theorem}\label{thm:main-ecr}
      Let $F$ be a non-archimedean local field of characteristic zero.
      Let $\psi$ be a non-tempered A-parameter of $\rG_2(F)$ and $\psi' = \iota \circ \psi$ be the lift of $\psi$ to $\PGSO_8(F)$.
      Then, the following statements hold.
      \begin{enumerate}
            \item
            We have 
            \begin{align*}
                  \Pi_{\widehat{\psi}} &= \widehat{\Pi_{\psi}}, \\
                  \Pi_{\widehat{\psi'}} &= \widehat{\Pi_{\psi'}}.
            \end{align*}

            \item
            The distribution
            \begin{align*}
                  S\Theta_{\psi} = \sum_{\rho \in \Irr_{\C}(\cS_{\psi})} \tr\rho(s_{\psi}) \Theta_{\pi_{\psi}(\rho)}
            \end{align*}
            is stable.

            \item 
            The twisted endoscopic character identities hold:
            \begin{align*}
                  \sum_{\pi' \in \Pi_{\psi'}} \langle s s_{\psi'}, \pi' \rangle_{\psi'} \Theta_{\widetilde{\pi'}}
                  =
                  \Trans^{\widetilde{\PGSO_8}, *}_{H_s} S\Theta^{H_s}_{\psi_s},
            \end{align*}
            where $s \in \cS_{\psi'}$, $H_s$ is the endoscopic group of $\PGSO_8 \rtimes \theta$ associated with $s \rtimes \theta$, $\psi_s$ is the A-parameter of $H_s$ associated with $\psi'$, and $S\Theta^{H_s}_{\psi_s}$ is the stable distribution of $H_s$ associated with $\psi_s$.

            \item 
            The endoscopic character identities for $A$-parameters $\psi$ of $\rG_2(F)$ hold:

            \begin{align*}
                  \Theta_{\psi, s}
                  =
                  \Trans^{\rG_2, *}_{H_s} S\Theta^{H_s}_{\psi_s},
            \end{align*}
            where $s \in \cS_{\psi}$, $H_s$ is the endoscopic group of $\rG_2(F)$ associated with $s$, $\psi_s$ is the A-parameter of $H_s$ associated with $\psi$, and $S\Theta^{H_s}_{\psi_s}$ is the stable distribution of $H_s$ associated with $\psi_s$.
      \end{enumerate}
\end{theorem}

\begin{proof}
      These follow from Propositions \ref{prop:principal-ecr}, \ref{prop:subregular-trivial-ecr}, \ref{prop:subregular-quadratic-ecr}, \ref{prop:subregular-cubic-ecr}, \ref{prop:subregular-full-ecr}, \ref{prop:co-tempered-long-short-ecr}, \ref{prop:mixed-long-root-ecr}, \ref{prop:mixed-short-root-trivial-ecr}, \ref{prop:mixed-short-root-quadratic-ecr} and Proposition \ref{prop:co-tempered-long-short-unbounded-ecr}.
\end{proof}



\bibliographystyle{plain}
\bibliography{bibliography}
\end{document}
